% ROZDZIAL 8 — Iteratory, ranges i views (C++20)
% Fragment do włączenia po rozdziale 7.
% Wymagane pakiety: listings, array, amsmath (opcjonalnie), hyperref.
% Użycie: % ROZDZIAL 8 — Iteratory, ranges i views (C++20)
% Fragment do włączenia po rozdziale 7.
% Wymagane pakiety: listings, array, amsmath (opcjonalnie), hyperref.
% Użycie: % ROZDZIAL 8 — Iteratory, ranges i views (C++20)
% Fragment do włączenia po rozdziale 7.
% Wymagane pakiety: listings, array, amsmath (opcjonalnie), hyperref.
% Użycie: % ROZDZIAL 8 — Iteratory, ranges i views (C++20)
% Fragment do włączenia po rozdziale 7.
% Wymagane pakiety: listings, array, amsmath (opcjonalnie), hyperref.
% Użycie: \input{rozdzial_8_iteratory_ranges_views_cpp.tex}

\chapter[Iteratory, Ranges i Views]{Iteratory, zakresy i widoki: \textit{Ranges} oraz \textit{Views} w C++20}
\label{chap:ranges-views-cpp}

\section{Dlaczego wprowadzono Ranges?}

Klasyczne algorytmy STL operują na \emph{parze iteratorów}: pierwszym
wskazującym początek danych i drugim wyznaczającym koniec. Dzięki temu
\texttt{std::sort} lub \texttt{std::find} mogą współpracować z wieloma
kontenerami. Wymaga to jednak przekazywania obu granic i dodatkowej
ostrożności, aby pochodziły z tego samego zakresu.

\begin{lstlisting}
#include <algorithm>
#include <vector>

std::vector<int> xs{4, 1, 3, 2};
std::sort(xs.begin(), xs.end()); // STL before C++20
std::ranges::sort(xs);          // C++20
\end{lstlisting}

Nowy model pozwala przekazać \emph{zakres} (\textit{range}) jako jeden
obiekt oraz formułować wymagania wobec niego poprzez \emph{koncepcje}
(\textit{concepts}, rozdział~5). Oddzielną innowacją są \emph{widoki}
(\textit{views}): lekkie reprezentacje danych, często przekształcane
\textbf{leniwe}, to znaczy dopiero podczas ich przeglądania.

\begin{lstlisting}
#include <ranges>

std::vector<int> xs{1, 2, 3, 4, 5, 6};
auto pipeline = xs
    | std::views::filter([](int x) { return x % 2 == 0; })
    | std::views::transform([](int x) { return x * x; });

// No vector of squares has been created.
// Iteration produces 4, 16, 36.
for (int x : pipeline) { /* use x */ }
\end{lstlisting}

\paragraph{Założenia.}
W całym rozdziale przyjmujemy standard \textbf{C++20} oraz nagłówki
\texttt{<algorithm>}, \texttt{<iterator>}, \texttt{<ranges>},
\texttt{<concepts>}, \texttt{<vector>} i pozostałe nagłówki wymienione
w przykładach. W dalszych sekcjach posłużymy się semantyką przenoszenia
(rozdział~6) i lambdami (rozdział~7).

\section{Iteratory: uogólnienie wskaźników}
\label{sec:cpp8-iterators}

\subsection{Operacje \texttt{begin}, \texttt{end}, \texttt{*} i \texttt{++}}

Iterator jest obiektem umożliwiającym poruszanie się po elementach
zakresu. Wskaźnik \texttt{int*} jest przykładem iteratora, ale iterator
może także przechowywać indeks, wskaźnik do węzła listy lub stan
wewnętrzny generatora. Podstawowe operacje przedstawia program:

\begin{lstlisting}
#include <vector>

std::vector<int> v{10, 20, 30};
auto it = v.begin();     // points to 10
int a = *it;            // 10
++it;                   // points to 20
int b = *it;            // 20
bool ok = it != v.end();
\end{lstlisting}

Iterator \texttt{end()} zwykle oznacza pozycję \emph{za ostatnim}
elementem; nie wolno go dereferencjonować. Dzięki temu pusty zakres
spełnia \texttt{begin()==end()}, a przechodzenie po elementach ma
jednolity schemat:

\begin{lstlisting}
for (auto it = v.begin(); it != v.end(); ++it) {
    // *it is a valid element
}
\end{lstlisting}

\paragraph{Pętla zakresowa.}
Zapis \texttt{for (auto\& x : v)} jest wygodniejszą postacią iteracji,
ale w dalszym ciągu wykorzystuje operacje początku, końca,
porównania, zwiększania i dereferencji. Aby \emph{modyfikować} elementy,
stosujemy \texttt{auto\&}; aby uniknąć kopiowania bez modyfikowania,
\texttt{const auto\&}; dla niewielkich wartości liczbowych często
wystarczy \texttt{auto}.

\subsection{Iteratory stałe i nieważność iteratorów}

\begin{lstlisting}
std::vector<int> v{1, 2, 3};
auto m = v.begin();   // can modify *m
auto c = v.cbegin();  // cannot modify *c
*m = 10;
// *c = 20;         // compilation error
\end{lstlisting}

\texttt{const\_iterator} chroni elementy przed modyfikacją poprzez
ten iterator; nie oznacza, że cały kontener jest niezmienny. Równie
ważna jest \emph{nieważność iteratorów} (\textit{iterator invalidation}).
Na przykład ponowna alokacja pamięci przez \texttt{std::vector::push\_back}
może unieważnić \textbf{wszystkie} iteratory, wskaźniki i referencje do
jego elementów.

\begin{lstlisting}
std::vector<int> v{1, 2, 3};
auto it = v.begin();
v.push_back(4);       // may reallocate storage
// *it;              // unsafe if reallocation happened
it = v.begin();      // obtain a fresh iterator
\end{lstlisting}

Reguły zależą od kontenera i operacji. W \texttt{std::list} wstawienie
elementu nie unieważnia iteratorów do pozostałych elementów, ale nie
oznacza to, że można dalej używać iteratora do elementu usuniętego.

\section{Kategorie iteratorów i ich wymagania}
\label{sec:cpp8-iterator-categories}

Zestaw koncepcji C++20 opisuje operacje i własności iteratorów, a nie
jedynie nazwy klas. Najważniejsze kategorie zestawia tabela.

\begin{center}
\small
\begin{tabular}{|p{5.2cm}|p{7.3cm}|}
\hline
\textbf{Koncepcja} & \textbf{Możliwości i przykłady} \\
\hline
\texttt{input\_iterator} & Odczytywanie kolejnych wartości; przeglądanie może być jednokrotne (np. strumień wejściowy). \\
\hline
\texttt{output\_iterator} & Zapisywanie kolejnych wartości do celu (np. \texttt{back\_insert\_iterator}). \\
\hline
\texttt{forward\_iterator} & Wielokrotne przejścia, porównywanie iteratorów, krok do przodu. \\
\hline
\texttt{bidirectional\_iterator} & Dodatkowo krok wstecz, czyli \texttt{--it}; np. \texttt{list}. \\
\hline
\texttt{random\_access\_iterator} & Indeksowanie i skoki \texttt{it+n} w stałym czasie; np. \texttt{deque}. \\
\hline
\texttt{contiguous\_iterator} & Elementy leżą kolejno w pamięci; np. iterator \texttt{vector<int>} i surowy wskaźnik. \\
\hline
\end{tabular}
\end{center}

Hierarchia wymagań dla iteratorów odczytujących jest w przybliżeniu
następująca (każda kolejna koncepcja jest silniejsza):

\begin{center}
\texttt{input $\Rightarrow$ forward $\Rightarrow$ bidirectional}
\\[2pt]
\texttt{$\Rightarrow$ random\_access $\Rightarrow$ contiguous}.
\end{center}

Koncepcja \texttt{output\_iterator} dotyczy osobnej rodziny wymagań.
Nowe koncepcje nie są w każdym szczególe identyczne ze starymi
kategoriami \emph{LegacyIterator}. Na przykład model
\texttt{forward\_iterator} nie wymaga zwracania referencji
przy dereferencji.

\begin{lstlisting}
#include <concepts>
#include <iterator>
#include <list>
#include <vector>

static_assert(std::contiguous_iterator<int*>);
static_assert(std::random_access_iterator<
    std::vector<int>::iterator>);
static_assert(std::bidirectional_iterator<
    std::list<int>::iterator>);
static_assert(!std::random_access_iterator<
    std::list<int>::iterator>);
\end{lstlisting}

\paragraph{Dlaczego to ważne?}
\texttt{std::ranges::sort} wymaga iteratorów o dostępie swobodnym,
więc można użyć go do \texttt{std::vector}, ale \textbf{nie} do
\texttt{std::list}. Listę sortuje się metodą \texttt{list::sort()}.
Wybór algorytmu jest zatem ograniczony kategorią iteratora.

\section{Własny iterator i własny zakres}
\label{sec:cpp8-own-iterator}

Zbudujemy iterator generujący kolejne liczby całkowite. Nie będzie on
przechowywać ani wskaźnika, ani kontenera: całą potrzebną informacją
jest bieżąca wartość. Zdefiniujemy \emph{iterator do przodu}
(\texttt{forward\_iterator}).

\begin{lstlisting}
#include <concepts>
#include <cstddef>
#include <iterator>
#include <ranges>

struct CountingIterator {
    using value_type = int;
    using difference_type = std::ptrdiff_t;
    using iterator_concept = std::forward_iterator_tag;

    int current = 0;

    int operator*() const { return current; }

    CountingIterator& operator++() {
        ++current;
        return *this;
    }
    CountingIterator operator++(int) {
        auto old = *this;
        ++*this;
        return old;
    }
    bool operator==(const CountingIterator&) const = default;
};

struct CountingRange {
    int first = 0;
    int last = 0;

    CountingIterator begin() const { return {first}; }
    CountingIterator end() const { return {last}; }
};

static_assert(std::forward_iterator<CountingIterator>);
static_assert(std::ranges::forward_range<CountingRange>);
\end{lstlisting}

Teraz kod \texttt{for (int i : CountingRange\{3, 7\})} wypisze
\texttt{3, 4, 5, 6}. Funkcje \texttt{begin()} i \texttt{end()}
wyznaczają \emph{półotwarty} zakres \([3,7)\). Prawidłowe użycie
zakłada \texttt{first <= last}; dla \texttt{first > last} nasz
minimalny iterator nie zbiegnie do wartości końcowej.

\paragraph{Co trzeba zagwarantować?}
Samo przejście \texttt{static\_assert} nie dowodzi wszystkich
semantycznych własności iteratora. Dla \texttt{forward\_iterator}
ważna jest m.in. własność \emph{multi-pass}: kopiowanie iteratora
i przejście jedną kopią nie powinno zmieniać wartości zwracanych
przez drugą. Nasz iterator ją spełnia, ponieważ każda kopia ma własne
pole \texttt{current}.

\section{Zwykłe \textit{range} i \textit{sentinel}}
\label{sec:cpp8-range-sentinel}

\emph{Zakres} jest obiektem, dla którego można otrzymać początek
\texttt{std::ranges::begin(r)} i koniec \texttt{std::ranges::end(r)}.
Typ wyniku obu wywołań \textbf{nie musi być identyczny}. Obiekt
wyznaczający koniec to \emph{sentinel}; musi dać się porównać z
iteratorem, ale sam nie musi obsługiwać jego pozostałych operacji.

\begin{lstlisting}
#include <ranges>
#include <vector>

std::vector<int> xs{1, 2, 3};
auto first = std::ranges::begin(xs);
auto last = std::ranges::end(xs);
while (first != last) { ++first; }
\end{lstlisting}

W tym przypadku obie granice mają ten sam typ: zakres spełnia
koncepcję \texttt{common\_range}. W innych zakresach koniec może
wyznaczać osobny obiekt sentinela. Odpowiednia koncepcja
\texttt{std::sentinel\_for} określa warunki porównywania granic:

\begin{lstlisting}
static_assert(std::ranges::common_range<decltype(xs)>);
static_assert(std::sentinel_for<
    decltype(last), decltype(first)>);
\end{lstlisting}

Przykładem osobnego typu końca jest \texttt{std::default\_sentinel\_t}.

\paragraph{Rozmiar a liczba kroków.}
Operacja \texttt{std::ranges::size(r)} jest dostępna dla zakresów
spełniających \texttt{sized\_range}. Nie każdy zakres ma tanią operację
wyznaczenia rozmiaru. Funkcja \texttt{distance(r)} z przestrzeni
\texttt{std::ranges} może wymagać
\emph{przejścia przez wszystkie elementy}. Dla niektórych
jednokrotnych źródeł danych przejście oznacza również ich zużycie.

\section{Koncepcje zakresów w praktyce}

Własności zakresów sprawdzamy za pomocą koncepcji z rozdziału~5.
Najważniejsze z nich to:
\begin{itemize}
\item \texttt{input\_range} i \texttt{forward\_range},
\item \texttt{bidirectional\_range} i \texttt{random\_access\_range},
\item \texttt{contiguous\_range}, \texttt{common\_range} i \texttt{sized\_range}.
\end{itemize}

\begin{lstlisting}
#include <ranges>
#include <vector>

template <std::ranges::input_range R>
long long sum_integers(R&& range) {
    long long result = 0;
    for (auto value : range) result += value;
    return result;
}

std::vector<int> v{1, 2, 3};
long long s = sum_integers(v); // 6
\end{lstlisting}

Parametr \texttt{R\&\&} jest tutaj \emph{referencją przekazującą}
(\textit{forwarding reference}), a nie żądaniem przejęcia własności.
Funkcja obsługuje także źródła jednokrotne, gdyż przechodzi zakres tylko
raz. Dla naprawdę uniwersalnej funkcji trzeba osobno wyrazić warunki
na typ elementu i wykonywaną na nim operację dodawania.

\section{Algorytmy \texttt{std::ranges} i projekcje}
\label{sec:cpp8-algorithms}

Funkcje z rodziny \texttt{std::ranges}, np. \texttt{sort},
\texttt{find}, \texttt{copy} i \texttt{transform}, przyjmują
bezpośrednio zakres albo parę iterator--sentinel.
Ich wykonanie jest na ogół \emph{natychmiastowe} (w odróżnieniu od
samych widoków \texttt{views::transform}).

\begin{lstlisting}
#include <algorithm>
#include <iterator>
#include <vector>

std::vector<int> v{4, 2, 3, 1};
std::ranges::sort(v);  // modifies v immediately

auto pos = std::ranges::find(v, 3);
if (pos != v.end()) { *pos = 30; }

std::vector<int> out;
std::ranges::copy(v, std::back_inserter(out));
\end{lstlisting}

\subsection{Projekcja: sortowanie po wskazanym polu}

Algorytmy Ranges często posiadają argument \emph{projekcji}: funkcję
wywoływaną na elemencie przed porównaniem albo wyszukiwaniem. Dzięki
temu zamiast pisać własny komparator, można wskazać pole struktury.

\begin{lstlisting}
#include <string>
#include <vector>

struct Person {
    std::string name;
    int age;
};

std::vector<Person> people{
    {"Ada", 40}, {"Ewa", 25}, {"Jan", 32}
};

std::ranges::sort(people, {}, &Person::age);
// ascending ages: 25, 32, 40

auto it = std::ranges::find(people, 32, &Person::age);
// it->name == "Jan"
\end{lstlisting}

W wywołaniu \texttt{sort(people, \{\}, \&Person::age)} puste
nawiasy oznaczają domyślny komparator \texttt{std::ranges::less},
a wskaźnik do składowej to projekcja. Standard używa mechanizmu
\texttt{std::invoke}, omówionego w rozdziale~7.

\section{Widoki: \textit{view} nie jest synonimem kontenera}
\label{sec:cpp8-views}

\emph{Widok} jest zakresem przystosowanym do taniego tworzenia i
kopiowania lub przenoszenia \emph{samego obiektu widoku}. Nie musi mieć
własnych danych: może przechowywać referencję do kontenera, mały zestaw
iteratorów albo parametry generowania wartości.

\begin{lstlisting}
#include <ranges>
#include <vector>

std::vector<int> xs{10, 20, 30, 40};
auto first_two = xs | std::views::take(2);
// first_two refers to the original xs
\end{lstlisting}

Stworzenie widoku nie jest równoznaczne ze skopiowaniem elementów do
nowego kontenera. Odwołanie do danych źródłowych wymaga, by pozostawały
one ważne przez cały czas korzystania z widoku. Szczegółowe reguły
własności omówimy dalej.

\paragraph{Dwa znaczenia nazwy \texttt{views}.}
\texttt{std::views} jest aliasem przestrzeni nazw z obiektami
adaptującymi zakresy. \texttt{std::ranges::view} to natomiast
\emph{koncepcja} spełniana przez dany typ. Nie należy mylić ich
z pojęciem \texttt{std::ranges::range}.

\section{Leniwe potoki \texttt{filter}, \texttt{transform}, \texttt{take}}
\label{sec:cpp8-pipelines}

Najczęściej stosowane widoki łączymy za pomocą operatora \texttt{|}.
Adaptory przekazują sobie zakresy bez potrzeby budowania pośrednich
wektorów:

\begin{lstlisting}
std::vector<int> xs{1,2,3,4,5,6,7,8,9,10};

auto pipeline = xs
    | std::views::filter([](int x) {
          return x % 2 == 0;
      })
    | std::views::transform([](int x) {
          return x * x;
      })
    | std::views::take(3);

// Iteration produces 4, 16, 36.
\end{lstlisting}

\texttt{filter} pomija elementy niespełniające predykatu,
\texttt{transform} wyznacza nową wartość przy dereferencji,
a \texttt{take(3)} ogranicza liczbę wyników. Algorytm
\texttt{std::ranges::transform} działa od razu; natomiast
utworzenie widoku \texttt{views::transform} nie przelicza
wszystkich wyników.

\subsection{Kiedy następują obliczenia?}

\begin{lstlisting}
int calls = 0;
auto squares = xs | std::views::transform([&](int x) {
    ++calls;
    return x * x;
});
// calls == 0 here
int first = *squares.begin();
// calls has increased (transform was evaluated)
\end{lstlisting}

W tym przykładzie pokazujemy moment obliczenia, a nie zalecany styl
projektowania: \textbf{predykaty i transformacje powinny zasadniczo
unikać efektów ubocznych}. Ponowne przejście widoku, ponowna
dereferencja lub implementacyjnie dopuszczone buforowanie może
powodować inne liczby wywołań, niż sugerowałaby pojedyncza pętla.
Nie należy projektować logiki programu na podstawie rzekomej liczby
wywołań funkcji przez widok.

\paragraph{Ważny szczegół \texttt{filter\_view}.}
Dla pewnych kategorii zakresu \texttt{filter\_view} może zapamiętywać
pierwszy znaleziony element. Ponadto modyfikowanie elementu tak,
aby przestał spełniać predykat filtrowania, narusza wymagania
\texttt{filter\_view} i może prowadzić do niezdefiniowanego
zachowania. Predykat powinien być stabilny wobec przebiegu iteracji.

\section{Pozostałe przydatne adaptory C++20}
\label{sec:cpp8-adaptors}

\subsection{\texttt{iota}: wartości wyznaczane bez tablicy}

\begin{lstlisting}
auto ids = std::views::iota(1, 5);
// Iteration produces 1, 2, 3, 4.

// This is an unbounded range, limited by take(5):
auto first_five = std::views::iota(1)
                | std::views::take(5);
// Produces 1, 2, 3, 4, 5.
\end{lstlisting}

\texttt{iota} przydaje się do indeksowania lub generowania sekwencji.
Dla nieograniczonego widoku trzeba pamiętać o warunku zakończenia;
ponadto operacje arytmetyczne nie mogą przepełniać typu całkowitego
ze znakiem.

\subsection{\texttt{drop}, \texttt{take} i \texttt{reverse}}

\begin{lstlisting}
std::vector<int> data{1, 2, 3, 4, 5};
auto middle = data | std::views::drop(1)
                   | std::views::take(3);
// 2, 3, 4

auto backwards = data | std::views::reverse;
// 5, 4, 3, 2, 1
\end{lstlisting}

\texttt{reverse} wymaga co najmniej zakresu dwukierunkowego;
\texttt{take} i \texttt{drop} są przydatne także dla sekwencji
generowanych leniwie. Dla niektórych widoków koszt \texttt{begin()}
lub dojścia do pierwszego wyniku może być większy niż stały.

\subsection{\texttt{keys}, \texttt{values} i \texttt{split}}

\begin{lstlisting}
std::vector<std::pair<int, std::string>> pairs{
    {1, "one"}, {2, "two"}
};
for (int key : pairs | std::views::keys) { /* 1, 2 */ }
for (const auto& value : pairs | std::views::values) {
    /* "one", "two" */
}

std::string csv = "aa,b,cc";
for (auto part : csv | std::views::split(',')) {
    std::string piece(part.begin(), part.end());
    // Copy each part into an owning string if needed.
}
\end{lstlisting}

\texttt{split} nie tworzy automatycznie \texttt{std::string} dla
każdego fragmentu; otrzymujemy zakresy wskazujące fragmenty źródła.
Zachowując wynik na później, trzeba zapewnić odpowiedni czas życia
źródła. Nie należy też zakładać, że splitowanie tekstu CSV obsługuje
cudzysłowy i escapowanie pól zgodnie z pełnym formatem CSV.

\paragraph{Uwaga na wersję standardu.}
\texttt{std::views::zip} i \texttt{std::views::enumerate} są
konstrukcjami \textbf{C++23}, podobnie jak \texttt{std::ranges::to}.
Nie są dostępne w czystym C++20. W standardzie C++20 wynik można zebrać do kontenera np.
algorytmem \texttt{ranges::copy} z iteratorem wstawiającym
\texttt{std::back\_inserter}.

\section{Materializacja wyniku i koszt przetwarzania}
\label{sec:cpp8-materialize}

Jeśli wynik ma być niezależny od źródła i przechowywany na później,
możemy zamienić potok w kontener. W C++20 nie ma standardowego
\texttt{std::ranges::to<std::vector>()}.

\begin{lstlisting}
#include <algorithm>
#include <iterator>
#include <vector>

std::vector<int> data{1,2,3,4,5,6};
auto even = data | std::views::filter([](int x) {
    return x % 2 == 0;
});

std::vector<int> result;
std::ranges::copy(even, std::back_inserter(result));
// result == {2, 4, 6}
\end{lstlisting}

Taka materializacja wymaga pamięci proporcjonalnej do wielkości wyniku.
Zaletą potoku jest unikanie pośrednich alokacji i często możliwość
\emph{łączenia} transformacji w jednym przebiegu. Nie jest to jednak
gwarancja, że widoki zawsze będą szybsze: koszty iteracji, predykatów,
przepływu sterowania i optymalizacji zależą od konkretnego kodu i
kompilatora. W zadaniach krytycznych wykonuj pomiary.

\section{Własność danych, czas życia i \texttt{borrowed\_range}}
\label{sec:cpp8-lifetime}

To najważniejsza część praktycznego korzystania z Ranges. Widoki
\textbf{nie zawsze posiadają} bazowe dane. W przypadku
\texttt{std::views::all} przeprowadzamy istotne rozróżnienie:

\begin{lstlisting}
std::vector<int> source{2, 4, 6};
auto borrowed = std::views::all(source);
// Refers to source: source must stay alive.

auto owned = std::views::all(std::vector<int>{3, 5, 7});
// In C++20, this can be an owning_view.
// The temporary vector is moved into the view.
\end{lstlisting}

Nie oznacza to, że \emph{każdy} widok utworzony z wartości tymczasowej
przejmuje własność, ani że wszystkie widoki zachowują się tak samo.
Różne adaptory przechowują różne rodzaje obiektów bazowych.

\subsection{Wskaźniki wiszące wewnątrz potoku}

\begin{lstlisting}
auto make_bad_view() {
    std::vector<int> local{1, 2, 3};
    return local | std::views::filter([](int x) {
        return x > 1;
    });
}
// Using the returned view would access destroyed local.
\end{lstlisting}

Ponieważ \texttt{local} jest \emph{lvalue}, adapter przechowuje
nieposiadający danych widok do lokalnego kontenera. Po wyjściu z
funkcji \texttt{local} zostaje zniszczony, a zwrócony widok zawiera
nieaktualne odwołanie. Bezpieczne wyjścia to utworzenie widoku z
przekazanego na własność kontenera (np. obiektu tymczasowego lub
\texttt{std::move(local)}) albo zwrócenie zwykłego kontenera wynikowego.

\begin{lstlisting}
auto make_owning_view() {
    std::vector<int> local{1, 2, 3};
    return std::views::all(std::move(local))
         | std::views::filter([](int x) { return x > 1; });
}
// The returned view owns its vector through owning_view.
\end{lstlisting}

\subsection{Co naprawdę mówi \texttt{borrowed\_range}?}

Koncepcja \texttt{borrowed\_range<R>} oznacza w uproszczeniu, że
iteratory otrzymane z zakresu typu \texttt{R} mogą pozostać ważne
po zniszczeniu \emph{samego obiektu zakresu} (o ile dane, do których
odsyłają, nadal istnieją). Nie jest to synonim ``widok nie posiada
danych'' ani ``wszystkie wskaźniki są zawsze bezpieczne''.

\begin{lstlisting}
#include <concepts>
#include <ranges>
#include <vector>

static_assert(std::ranges::borrowed_range<
    std::ranges::ref_view<std::vector<int>>>);
static_assert(!std::ranges::borrowed_range<
    std::vector<int>>);
\end{lstlisting}

\subsection{\texttt{std::ranges::dangling} jako zabezpieczenie}

Algorytm zwracający iterator do \emph{tymczasowego} kontenera nie może
bezpiecznie zwrócić rzeczywistego iteratora do zniszczonych danych.
Dlatego pewne algorytmy Ranges zwracają specjalny typ
\texttt{std::ranges::dangling}.

\begin{lstlisting}
#include <type_traits>

using Result = decltype(std::ranges::find(
    std::vector<int>{1, 2, 3}, 2));
static_assert(std::same_as<Result, std::ranges::dangling>);
\end{lstlisting}

Nie należy jednak oczekiwać, że mechanizm ten wykryje wszystkie
problemy z czasem życia. Jeżeli zapisujemy nieposiadający danych
widok do lokalnej zmiennej, kompilator może nie zauważyć, że docelowy
kontener umrze wcześniej niż widok.

\section{\texttt{std::span}, \texttt{std::string\_view} i \texttt{subrange}}
\label{sec:cpp8-span}

\texttt{std::span<T>} opisuje ciągły fragment pamięci jako wskaźnik
i rozmiar (lub rozmiar znany w czasie kompilacji). Jest to lekka,
\textbf{nieposiadająca danych} alternatywa dla pary
\texttt{(T*, size)}.

\begin{lstlisting}
#include <span>

void normalize(std::span<double> xs) {
    for (double& x : xs) x /= 100.0;
}

std::vector<double> v{100, 200, 300};
normalize(v);  // v becomes {1, 2, 3}
\end{lstlisting}

\texttt{std::span<const T>} umożliwia odczyt bez modyfikowania
elementów, a \texttt{std::span<T>} pozwala modyfikować elementy, mimo
że sam widok nie posiada pamięci.

\begin{lstlisting}
int a[]{1, 2, 3, 4};
std::span<int> all{a};
auto head = all.first(2);
head[1] = 20; // a[1] == 20
\end{lstlisting}

\texttt{std::string\_view} spełnia podobną rolę dla tekstu; nie
alokuje pamięci na znaki (oraz nie musi kończyć się bajtem zero).
Trzeba uważać, aby nie zapisać go po zniszczeniu źródłowego
\texttt{std::string}.

\begin{lstlisting}
#include <string_view>

std::string text = "abcdef";
std::string_view view{text};
std::string_view prefix = view.substr(0, 3); // "abc"
// Do not use view after text is destroyed or invalidated.
\end{lstlisting}

\texttt{std::ranges::subrange<I, S>} grupuje iterator \texttt{I}
i sentinel \texttt{S} w zakres bez kopiowania elementów. Także tutaj
obowiązuje odpowiedzialność za czas życia wskazywanych danych.

\begin{lstlisting}
auto sub = std::ranges::subrange(v.begin(), v.end());
// sub refers to elements of v; it does not own v.
\end{lstlisting}

\paragraph{Ostrożnie z realokacją.}
\texttt{std::span} do elementów wektora, \texttt{string\_view} do
znaków łańcucha oraz \texttt{subrange} z iteratorów do kontenera mogą
stać się niepoprawne po operacji zmieniającej miejsce przechowywania
danych. Przenoszenie lub realokacja nie są automatycznie śledzone.

\section{Własny widok za pomocą \texttt{view\_interface}}
\label{sec:cpp8-custom-view}

Na końcu zbudujemy \emph{własny widok} generujący kwadraty liczb
całkowitych od zera do \(n-1\). Użyjemy
\texttt{std::ranges::view\_interface}, które dostarcza wygodnych
operacji opartych na \texttt{begin()} i \texttt{end()}. W klasie
pozostają wyłącznie parametry generatora, więc nie ma alokacji pamięci
na elementy sekwencji.

\begin{lstlisting}
#include <concepts>
#include <cstddef>
#include <iterator>
#include <ranges>

struct SquaresView : std::ranges::view_interface<SquaresView> {
    struct iterator {
        using value_type = int;
        using difference_type = std::ptrdiff_t;
        using iterator_concept = std::forward_iterator_tag;

        int index = 0;
        int operator*() const { return index * index; }
        iterator& operator++() {
            ++index;
            return *this;
        }
        iterator operator++(int) {
            auto old = *this;
            ++*this;
            return old;
        }
        bool operator==(const iterator&) const = default;
    };

    int n = 0;
    SquaresView() = default;
    explicit SquaresView(int count) : n(count) {}

    iterator begin() const { return {0}; }
    iterator end() const { return {n}; }
};

static_assert(std::ranges::view<SquaresView>);
static_assert(std::ranges::forward_range<SquaresView>);
\end{lstlisting}

\begin{lstlisting}
SquaresView squares{6};
for (int x : squares) { /* 0, 1, 4, 9, 16, 25 */ }

// Our view works with standard adapters:
auto even_squares = squares | std::views::filter(
    [](int x) { return x % 2 == 0; });
// Iteration yields 0, 4, 16.
\end{lstlisting}

Konstruktor z argumentem \texttt{count} powinien otrzymać wartość
nieujemną. Ponieważ liczymy \texttt{index * index} w typie
\texttt{int}, należy ponadto ograniczyć rozmiar zakresu, aby uniknąć
przepełnienia. W rozwiązaniu produkcyjnym należałoby te warunki
weryfikować oraz dopracować obsługę szerszych typów całkowitych.

\paragraph{Widok a adapter.}
\texttt{SquaresView} jest własnym \emph{widokiem}, który można
przekazać do standardowych adapterów. Nie jest jeszcze własnym
\emph{range adaptor closure}, czyli obiektem pozwalającym zapisać
\texttt{some\_range | squares\_adaptor}. Implementacja takiego
adaptera stanowi osobny poziom abstrakcji i nie jest konieczna do
budowania własnych widoków.

\section{Różnica między \texttt{transform} i \texttt{views::transform}}

Poniższe operacje mogą prowadzić do takich samych liczb, ale różnią
się \emph{czasem wykonania} i \emph{sposobem przechowywania wyniku}:

\begin{lstlisting}
std::vector<int> source{1, 2, 3};
std::vector<int> eager;
std::ranges::transform(source,
    std::back_inserter(eager),
    [](int x) { return x * x; });
// eager owns {1, 4, 9} right now

auto lazy = source | std::views::transform(
    [](int x) { return x * x; });
// lazy does not own a vector of squared values
\end{lstlisting}

Wariant pierwszy wymaga pamięci na trzy elementy i kończy obliczenia
przed kolejną instrukcją. Wariant drugi może obliczać wartości dopiero
przy dereferencji iteratorów. Po zmianie \texttt{source} kolejne
odczyty widoku \texttt{lazy} mogą odzwierciedlać nowe dane.

\section{Projektowanie API z wykorzystaniem Ranges}
\label{sec:cpp8-api}

Jeżeli funkcja tylko przegląda elementy dowolnego źródła, preferuj
parametr ograniczony koncepcją \texttt{input\_range}. Dla wielokrotnych
przejść korzystaj z \texttt{forward\_range}. Jeśli potrzebujesz
ciągłego bloku pamięci i znasz typ elementów, wygodnym wyborem jest
\texttt{std::span<T>}.

\begin{lstlisting}
template <std::ranges::input_range R>
void print_each(R&& r) {
    for (auto&& x : r) {
        // Print or consume x.
    }
}

void process_buffer(std::span<const double> values) {
    for (double value : values) {
        // Read numeric values.
    }
}
\end{lstlisting}

Zależnie od celu można zwracać widok albo kontener. Widok jest
wygodny, gdy chcemy dalej budować potok i mamy kontrolę nad czasem
życia źródła. Kontener jest często bezpieczniejszy, gdy zwracany
wynik ma istnieć niezależnie od wszystkich argumentów.

\section{Najczęstsze błędy i praktyczne reguły}
\label{sec:cpp8-pitfalls}

\begin{enumerate}
\item \textbf{Dereferencja \texttt{end()}.} Iterator końcowy jest
granicą, a nie elementem. Przed \texttt{*it} trzeba mieć
\texttt{it != end}.

\item \textbf{Mylenie kategorii iteratorów.} Lista nie obsługuje
skoków w stałym czasie. Nie można użyć dla niej algorytmu
\texttt{ranges::sort}; należy wywołać metodę \texttt{list::sort}.

\item \textbf{Traktowanie widoku jak kontenera.} \texttt{views::filter}
nie tworzy kopii wybranych elementów. Zmiany i zniszczenie źródła
mogą wpływać na wynik.

\item \textbf{Wiszące referencje w lambdach.} Zwrócony widok
z predykatem przechwytującym lokalny obiekt przez referencję może
być niebezpieczny nawet wtedy, gdy sam kontener ma odpowiednio długi
czas życia.

\item \textbf{Efekty uboczne w predykacie.} Liczba i kolejność
wywołań w leniwym potoku nie są odpowiednią podstawą działania
programu; predykaty filtrowania powinny zachować stabilną semantykę.

\item \textbf{Zakładanie stałego kosztu operacji.} Obliczenie
\texttt{distance} lub pierwszego pasującego elementu może wymagać
przejrzenia dużej części źródła.

\item \textbf{Mieszanie standardów.} \texttt{ranges::to},
\texttt{views::zip} i \texttt{views::enumerate} nie należą do C++20.

\item \textbf{Używanie nieważnych iteratorów.} Realokacja wektora
lub usuwanie elementu może unieważnić iteratory także wtedy, gdy
znajdują się wewnątrz wcześniej utworzonego widoku.
\end{enumerate}

\section{Której konstrukcji użyć?}

\begin{center}
\small
\begin{tabular}{|p{6.0cm}|p{6.5cm}|}
\hline
\textbf{Potrzeba} & \textbf{Zalecany wybór} \\
\hline
Przejście po zwykłym kontenerze & Pętla \texttt{for (auto\& x : v)} \\
\hline
Jednorazowe sortowanie albo wyszukiwanie & \texttt{std::ranges::sort}, \texttt{std::ranges::find} \\
\hline
Leniwe filtrowanie i przekształcanie & \texttt{std::views::filter}, \texttt{std::views::transform} \\
\hline
Generowanie sekwencji indeksów & \texttt{std::views::iota} \\
\hline
Odczyt lub modyfikacja ciągłego bufora & \texttt{std::span<T>} (bez własności) \\
\hline
Lekki odczyt fragmentu tekstu & \texttt{std::string\_view} (bez własności) \\
\hline
Zakres wyznaczony iteratorami & \texttt{std::ranges::subrange} \\
\hline
Wynik niezależny od źródła & Kontener, np. \texttt{std::vector} \\
\hline
Własny generator wartości & Typ modelujący \texttt{std::ranges::view} \\
\hline
Wydajny ogólny algorytm generyczny & Szablon z \texttt{std::ranges::input\_range} lub silniejszą koncepcją \\
\hline
\end{tabular}
\end{center}

\paragraph{Pytania kontrolne.}
\begin{enumerate}
\item Dlaczego iterator \texttt{end()} nie jest elementem kolekcji?
\item Co odróżnia \texttt{forward\_iterator} od \texttt{input\_iterator}?
\item Czy \texttt{std::ranges::sort} działa dla \texttt{std::list}? Dlaczego?
\item Po co \texttt{std::ranges} dopuszcza sentinel innego typu niż iterator?
\item W którym momencie \texttt{views::transform} wykonuje funkcję?
\item Dlaczego \texttt{views::filter} nie powinien zmieniać elementów, od których zależy jego predykat?
\item Czym różni się \texttt{views::all(lvalue)} od \texttt{views::all(rvalue)} dla \texttt{std::vector}?
\item Co gwarantuje, a czego nie gwarantuje \texttt{borrowed\_range}?
\item Dlaczego \texttt{ranges::find(std::vector\{\dots\}, x)} może zwrócić \texttt{dangling}?
\item Jak zebrać elementy \texttt{filter\_view} do wektora w czystym C++20?
\item Czym różni się \texttt{std::span<const T>} od \texttt{std::vector<T>}?
\item Co należy zaimplementować, aby własny obiekt był \texttt{forward\_range}?
\end{enumerate}

\paragraph{Dalsza lektura.}
Dobrym punktem odniesienia jest dokumentacja \emph{cppreference}:
\emph{Ranges library}, \emph{Iterator library}, \emph{Range adaptors},
\emph{std::ranges::borrowed\_range}, \emph{std::ranges::dangling}
oraz \emph{std::ranges::view\_interface}. Warto testować kod z opcjami kompilatora
\texttt{-std=c++20} i \texttt{-Wall -Wextra -Wpedantic},
a także korzystać z AddressSanitizer przy badaniu czasu życia danych.


\chapter[Iteratory, Ranges i Views]{Iteratory, zakresy i widoki: \textit{Ranges} oraz \textit{Views} w C++20}
\label{chap:ranges-views-cpp}

\section{Dlaczego wprowadzono Ranges?}

Klasyczne algorytmy STL operują na \emph{parze iteratorów}: pierwszym
wskazującym początek danych i drugim wyznaczającym koniec. Dzięki temu
\texttt{std::sort} lub \texttt{std::find} mogą współpracować z wieloma
kontenerami. Wymaga to jednak przekazywania obu granic i dodatkowej
ostrożności, aby pochodziły z tego samego zakresu.

\begin{lstlisting}
#include <algorithm>
#include <vector>

std::vector<int> xs{4, 1, 3, 2};
std::sort(xs.begin(), xs.end()); // STL before C++20
std::ranges::sort(xs);          // C++20
\end{lstlisting}

Nowy model pozwala przekazać \emph{zakres} (\textit{range}) jako jeden
obiekt oraz formułować wymagania wobec niego poprzez \emph{koncepcje}
(\textit{concepts}, rozdział~5). Oddzielną innowacją są \emph{widoki}
(\textit{views}): lekkie reprezentacje danych, często przekształcane
\textbf{leniwe}, to znaczy dopiero podczas ich przeglądania.

\begin{lstlisting}
#include <ranges>

std::vector<int> xs{1, 2, 3, 4, 5, 6};
auto pipeline = xs
    | std::views::filter([](int x) { return x % 2 == 0; })
    | std::views::transform([](int x) { return x * x; });

// No vector of squares has been created.
// Iteration produces 4, 16, 36.
for (int x : pipeline) { /* use x */ }
\end{lstlisting}

\paragraph{Założenia.}
W całym rozdziale przyjmujemy standard \textbf{C++20} oraz nagłówki
\texttt{<algorithm>}, \texttt{<iterator>}, \texttt{<ranges>},
\texttt{<concepts>}, \texttt{<vector>} i pozostałe nagłówki wymienione
w przykładach. W dalszych sekcjach posłużymy się semantyką przenoszenia
(rozdział~6) i lambdami (rozdział~7).

\section{Iteratory: uogólnienie wskaźników}
\label{sec:cpp8-iterators}

\subsection{Operacje \texttt{begin}, \texttt{end}, \texttt{*} i \texttt{++}}

Iterator jest obiektem umożliwiającym poruszanie się po elementach
zakresu. Wskaźnik \texttt{int*} jest przykładem iteratora, ale iterator
może także przechowywać indeks, wskaźnik do węzła listy lub stan
wewnętrzny generatora. Podstawowe operacje przedstawia program:

\begin{lstlisting}
#include <vector>

std::vector<int> v{10, 20, 30};
auto it = v.begin();     // points to 10
int a = *it;            // 10
++it;                   // points to 20
int b = *it;            // 20
bool ok = it != v.end();
\end{lstlisting}

Iterator \texttt{end()} zwykle oznacza pozycję \emph{za ostatnim}
elementem; nie wolno go dereferencjonować. Dzięki temu pusty zakres
spełnia \texttt{begin()==end()}, a przechodzenie po elementach ma
jednolity schemat:

\begin{lstlisting}
for (auto it = v.begin(); it != v.end(); ++it) {
    // *it is a valid element
}
\end{lstlisting}

\paragraph{Pętla zakresowa.}
Zapis \texttt{for (auto\& x : v)} jest wygodniejszą postacią iteracji,
ale w dalszym ciągu wykorzystuje operacje początku, końca,
porównania, zwiększania i dereferencji. Aby \emph{modyfikować} elementy,
stosujemy \texttt{auto\&}; aby uniknąć kopiowania bez modyfikowania,
\texttt{const auto\&}; dla niewielkich wartości liczbowych często
wystarczy \texttt{auto}.

\subsection{Iteratory stałe i nieważność iteratorów}

\begin{lstlisting}
std::vector<int> v{1, 2, 3};
auto m = v.begin();   // can modify *m
auto c = v.cbegin();  // cannot modify *c
*m = 10;
// *c = 20;         // compilation error
\end{lstlisting}

\texttt{const\_iterator} chroni elementy przed modyfikacją poprzez
ten iterator; nie oznacza, że cały kontener jest niezmienny. Równie
ważna jest \emph{nieważność iteratorów} (\textit{iterator invalidation}).
Na przykład ponowna alokacja pamięci przez \texttt{std::vector::push\_back}
może unieważnić \textbf{wszystkie} iteratory, wskaźniki i referencje do
jego elementów.

\begin{lstlisting}
std::vector<int> v{1, 2, 3};
auto it = v.begin();
v.push_back(4);       // may reallocate storage
// *it;              // unsafe if reallocation happened
it = v.begin();      // obtain a fresh iterator
\end{lstlisting}

Reguły zależą od kontenera i operacji. W \texttt{std::list} wstawienie
elementu nie unieważnia iteratorów do pozostałych elementów, ale nie
oznacza to, że można dalej używać iteratora do elementu usuniętego.

\section{Kategorie iteratorów i ich wymagania}
\label{sec:cpp8-iterator-categories}

Zestaw koncepcji C++20 opisuje operacje i własności iteratorów, a nie
jedynie nazwy klas. Najważniejsze kategorie zestawia tabela.

\begin{center}
\small
\begin{tabular}{|p{5.2cm}|p{7.3cm}|}
\hline
\textbf{Koncepcja} & \textbf{Możliwości i przykłady} \\
\hline
\texttt{input\_iterator} & Odczytywanie kolejnych wartości; przeglądanie może być jednokrotne (np. strumień wejściowy). \\
\hline
\texttt{output\_iterator} & Zapisywanie kolejnych wartości do celu (np. \texttt{back\_insert\_iterator}). \\
\hline
\texttt{forward\_iterator} & Wielokrotne przejścia, porównywanie iteratorów, krok do przodu. \\
\hline
\texttt{bidirectional\_iterator} & Dodatkowo krok wstecz, czyli \texttt{--it}; np. \texttt{list}. \\
\hline
\texttt{random\_access\_iterator} & Indeksowanie i skoki \texttt{it+n} w stałym czasie; np. \texttt{deque}. \\
\hline
\texttt{contiguous\_iterator} & Elementy leżą kolejno w pamięci; np. iterator \texttt{vector<int>} i surowy wskaźnik. \\
\hline
\end{tabular}
\end{center}

Hierarchia wymagań dla iteratorów odczytujących jest w przybliżeniu
następująca (każda kolejna koncepcja jest silniejsza):

\begin{center}
\texttt{input $\Rightarrow$ forward $\Rightarrow$ bidirectional}
\\[2pt]
\texttt{$\Rightarrow$ random\_access $\Rightarrow$ contiguous}.
\end{center}

Koncepcja \texttt{output\_iterator} dotyczy osobnej rodziny wymagań.
Nowe koncepcje nie są w każdym szczególe identyczne ze starymi
kategoriami \emph{LegacyIterator}. Na przykład model
\texttt{forward\_iterator} nie wymaga zwracania referencji
przy dereferencji.

\begin{lstlisting}
#include <concepts>
#include <iterator>
#include <list>
#include <vector>

static_assert(std::contiguous_iterator<int*>);
static_assert(std::random_access_iterator<
    std::vector<int>::iterator>);
static_assert(std::bidirectional_iterator<
    std::list<int>::iterator>);
static_assert(!std::random_access_iterator<
    std::list<int>::iterator>);
\end{lstlisting}

\paragraph{Dlaczego to ważne?}
\texttt{std::ranges::sort} wymaga iteratorów o dostępie swobodnym,
więc można użyć go do \texttt{std::vector}, ale \textbf{nie} do
\texttt{std::list}. Listę sortuje się metodą \texttt{list::sort()}.
Wybór algorytmu jest zatem ograniczony kategorią iteratora.

\section{Własny iterator i własny zakres}
\label{sec:cpp8-own-iterator}

Zbudujemy iterator generujący kolejne liczby całkowite. Nie będzie on
przechowywać ani wskaźnika, ani kontenera: całą potrzebną informacją
jest bieżąca wartość. Zdefiniujemy \emph{iterator do przodu}
(\texttt{forward\_iterator}).

\begin{lstlisting}
#include <concepts>
#include <cstddef>
#include <iterator>
#include <ranges>

struct CountingIterator {
    using value_type = int;
    using difference_type = std::ptrdiff_t;
    using iterator_concept = std::forward_iterator_tag;

    int current = 0;

    int operator*() const { return current; }

    CountingIterator& operator++() {
        ++current;
        return *this;
    }
    CountingIterator operator++(int) {
        auto old = *this;
        ++*this;
        return old;
    }
    bool operator==(const CountingIterator&) const = default;
};

struct CountingRange {
    int first = 0;
    int last = 0;

    CountingIterator begin() const { return {first}; }
    CountingIterator end() const { return {last}; }
};

static_assert(std::forward_iterator<CountingIterator>);
static_assert(std::ranges::forward_range<CountingRange>);
\end{lstlisting}

Teraz kod \texttt{for (int i : CountingRange\{3, 7\})} wypisze
\texttt{3, 4, 5, 6}. Funkcje \texttt{begin()} i \texttt{end()}
wyznaczają \emph{półotwarty} zakres \([3,7)\). Prawidłowe użycie
zakłada \texttt{first <= last}; dla \texttt{first > last} nasz
minimalny iterator nie zbiegnie do wartości końcowej.

\paragraph{Co trzeba zagwarantować?}
Samo przejście \texttt{static\_assert} nie dowodzi wszystkich
semantycznych własności iteratora. Dla \texttt{forward\_iterator}
ważna jest m.in. własność \emph{multi-pass}: kopiowanie iteratora
i przejście jedną kopią nie powinno zmieniać wartości zwracanych
przez drugą. Nasz iterator ją spełnia, ponieważ każda kopia ma własne
pole \texttt{current}.

\section{Zwykłe \textit{range} i \textit{sentinel}}
\label{sec:cpp8-range-sentinel}

\emph{Zakres} jest obiektem, dla którego można otrzymać początek
\texttt{std::ranges::begin(r)} i koniec \texttt{std::ranges::end(r)}.
Typ wyniku obu wywołań \textbf{nie musi być identyczny}. Obiekt
wyznaczający koniec to \emph{sentinel}; musi dać się porównać z
iteratorem, ale sam nie musi obsługiwać jego pozostałych operacji.

\begin{lstlisting}
#include <ranges>
#include <vector>

std::vector<int> xs{1, 2, 3};
auto first = std::ranges::begin(xs);
auto last = std::ranges::end(xs);
while (first != last) { ++first; }
\end{lstlisting}

W tym przypadku obie granice mają ten sam typ: zakres spełnia
koncepcję \texttt{common\_range}. W innych zakresach koniec może
wyznaczać osobny obiekt sentinela. Odpowiednia koncepcja
\texttt{std::sentinel\_for} określa warunki porównywania granic:

\begin{lstlisting}
static_assert(std::ranges::common_range<decltype(xs)>);
static_assert(std::sentinel_for<
    decltype(last), decltype(first)>);
\end{lstlisting}

Przykładem osobnego typu końca jest \texttt{std::default\_sentinel\_t}.

\paragraph{Rozmiar a liczba kroków.}
Operacja \texttt{std::ranges::size(r)} jest dostępna dla zakresów
spełniających \texttt{sized\_range}. Nie każdy zakres ma tanią operację
wyznaczenia rozmiaru. Funkcja \texttt{distance(r)} z przestrzeni
\texttt{std::ranges} może wymagać
\emph{przejścia przez wszystkie elementy}. Dla niektórych
jednokrotnych źródeł danych przejście oznacza również ich zużycie.

\section{Koncepcje zakresów w praktyce}

Własności zakresów sprawdzamy za pomocą koncepcji z rozdziału~5.
Najważniejsze z nich to:
\begin{itemize}
\item \texttt{input\_range} i \texttt{forward\_range},
\item \texttt{bidirectional\_range} i \texttt{random\_access\_range},
\item \texttt{contiguous\_range}, \texttt{common\_range} i \texttt{sized\_range}.
\end{itemize}

\begin{lstlisting}
#include <ranges>
#include <vector>

template <std::ranges::input_range R>
long long sum_integers(R&& range) {
    long long result = 0;
    for (auto value : range) result += value;
    return result;
}

std::vector<int> v{1, 2, 3};
long long s = sum_integers(v); // 6
\end{lstlisting}

Parametr \texttt{R\&\&} jest tutaj \emph{referencją przekazującą}
(\textit{forwarding reference}), a nie żądaniem przejęcia własności.
Funkcja obsługuje także źródła jednokrotne, gdyż przechodzi zakres tylko
raz. Dla naprawdę uniwersalnej funkcji trzeba osobno wyrazić warunki
na typ elementu i wykonywaną na nim operację dodawania.

\section{Algorytmy \texttt{std::ranges} i projekcje}
\label{sec:cpp8-algorithms}

Funkcje z rodziny \texttt{std::ranges}, np. \texttt{sort},
\texttt{find}, \texttt{copy} i \texttt{transform}, przyjmują
bezpośrednio zakres albo parę iterator--sentinel.
Ich wykonanie jest na ogół \emph{natychmiastowe} (w odróżnieniu od
samych widoków \texttt{views::transform}).

\begin{lstlisting}
#include <algorithm>
#include <iterator>
#include <vector>

std::vector<int> v{4, 2, 3, 1};
std::ranges::sort(v);  // modifies v immediately

auto pos = std::ranges::find(v, 3);
if (pos != v.end()) { *pos = 30; }

std::vector<int> out;
std::ranges::copy(v, std::back_inserter(out));
\end{lstlisting}

\subsection{Projekcja: sortowanie po wskazanym polu}

Algorytmy Ranges często posiadają argument \emph{projekcji}: funkcję
wywoływaną na elemencie przed porównaniem albo wyszukiwaniem. Dzięki
temu zamiast pisać własny komparator, można wskazać pole struktury.

\begin{lstlisting}
#include <string>
#include <vector>

struct Person {
    std::string name;
    int age;
};

std::vector<Person> people{
    {"Ada", 40}, {"Ewa", 25}, {"Jan", 32}
};

std::ranges::sort(people, {}, &Person::age);
// ascending ages: 25, 32, 40

auto it = std::ranges::find(people, 32, &Person::age);
// it->name == "Jan"
\end{lstlisting}

W wywołaniu \texttt{sort(people, \{\}, \&Person::age)} puste
nawiasy oznaczają domyślny komparator \texttt{std::ranges::less},
a wskaźnik do składowej to projekcja. Standard używa mechanizmu
\texttt{std::invoke}, omówionego w rozdziale~7.

\section{Widoki: \textit{view} nie jest synonimem kontenera}
\label{sec:cpp8-views}

\emph{Widok} jest zakresem przystosowanym do taniego tworzenia i
kopiowania lub przenoszenia \emph{samego obiektu widoku}. Nie musi mieć
własnych danych: może przechowywać referencję do kontenera, mały zestaw
iteratorów albo parametry generowania wartości.

\begin{lstlisting}
#include <ranges>
#include <vector>

std::vector<int> xs{10, 20, 30, 40};
auto first_two = xs | std::views::take(2);
// first_two refers to the original xs
\end{lstlisting}

Stworzenie widoku nie jest równoznaczne ze skopiowaniem elementów do
nowego kontenera. Odwołanie do danych źródłowych wymaga, by pozostawały
one ważne przez cały czas korzystania z widoku. Szczegółowe reguły
własności omówimy dalej.

\paragraph{Dwa znaczenia nazwy \texttt{views}.}
\texttt{std::views} jest aliasem przestrzeni nazw z obiektami
adaptującymi zakresy. \texttt{std::ranges::view} to natomiast
\emph{koncepcja} spełniana przez dany typ. Nie należy mylić ich
z pojęciem \texttt{std::ranges::range}.

\section{Leniwe potoki \texttt{filter}, \texttt{transform}, \texttt{take}}
\label{sec:cpp8-pipelines}

Najczęściej stosowane widoki łączymy za pomocą operatora \texttt{|}.
Adaptory przekazują sobie zakresy bez potrzeby budowania pośrednich
wektorów:

\begin{lstlisting}
std::vector<int> xs{1,2,3,4,5,6,7,8,9,10};

auto pipeline = xs
    | std::views::filter([](int x) {
          return x % 2 == 0;
      })
    | std::views::transform([](int x) {
          return x * x;
      })
    | std::views::take(3);

// Iteration produces 4, 16, 36.
\end{lstlisting}

\texttt{filter} pomija elementy niespełniające predykatu,
\texttt{transform} wyznacza nową wartość przy dereferencji,
a \texttt{take(3)} ogranicza liczbę wyników. Algorytm
\texttt{std::ranges::transform} działa od razu; natomiast
utworzenie widoku \texttt{views::transform} nie przelicza
wszystkich wyników.

\subsection{Kiedy następują obliczenia?}

\begin{lstlisting}
int calls = 0;
auto squares = xs | std::views::transform([&](int x) {
    ++calls;
    return x * x;
});
// calls == 0 here
int first = *squares.begin();
// calls has increased (transform was evaluated)
\end{lstlisting}

W tym przykładzie pokazujemy moment obliczenia, a nie zalecany styl
projektowania: \textbf{predykaty i transformacje powinny zasadniczo
unikać efektów ubocznych}. Ponowne przejście widoku, ponowna
dereferencja lub implementacyjnie dopuszczone buforowanie może
powodować inne liczby wywołań, niż sugerowałaby pojedyncza pętla.
Nie należy projektować logiki programu na podstawie rzekomej liczby
wywołań funkcji przez widok.

\paragraph{Ważny szczegół \texttt{filter\_view}.}
Dla pewnych kategorii zakresu \texttt{filter\_view} może zapamiętywać
pierwszy znaleziony element. Ponadto modyfikowanie elementu tak,
aby przestał spełniać predykat filtrowania, narusza wymagania
\texttt{filter\_view} i może prowadzić do niezdefiniowanego
zachowania. Predykat powinien być stabilny wobec przebiegu iteracji.

\section{Pozostałe przydatne adaptory C++20}
\label{sec:cpp8-adaptors}

\subsection{\texttt{iota}: wartości wyznaczane bez tablicy}

\begin{lstlisting}
auto ids = std::views::iota(1, 5);
// Iteration produces 1, 2, 3, 4.

// This is an unbounded range, limited by take(5):
auto first_five = std::views::iota(1)
                | std::views::take(5);
// Produces 1, 2, 3, 4, 5.
\end{lstlisting}

\texttt{iota} przydaje się do indeksowania lub generowania sekwencji.
Dla nieograniczonego widoku trzeba pamiętać o warunku zakończenia;
ponadto operacje arytmetyczne nie mogą przepełniać typu całkowitego
ze znakiem.

\subsection{\texttt{drop}, \texttt{take} i \texttt{reverse}}

\begin{lstlisting}
std::vector<int> data{1, 2, 3, 4, 5};
auto middle = data | std::views::drop(1)
                   | std::views::take(3);
// 2, 3, 4

auto backwards = data | std::views::reverse;
// 5, 4, 3, 2, 1
\end{lstlisting}

\texttt{reverse} wymaga co najmniej zakresu dwukierunkowego;
\texttt{take} i \texttt{drop} są przydatne także dla sekwencji
generowanych leniwie. Dla niektórych widoków koszt \texttt{begin()}
lub dojścia do pierwszego wyniku może być większy niż stały.

\subsection{\texttt{keys}, \texttt{values} i \texttt{split}}

\begin{lstlisting}
std::vector<std::pair<int, std::string>> pairs{
    {1, "one"}, {2, "two"}
};
for (int key : pairs | std::views::keys) { /* 1, 2 */ }
for (const auto& value : pairs | std::views::values) {
    /* "one", "two" */
}

std::string csv = "aa,b,cc";
for (auto part : csv | std::views::split(',')) {
    std::string piece(part.begin(), part.end());
    // Copy each part into an owning string if needed.
}
\end{lstlisting}

\texttt{split} nie tworzy automatycznie \texttt{std::string} dla
każdego fragmentu; otrzymujemy zakresy wskazujące fragmenty źródła.
Zachowując wynik na później, trzeba zapewnić odpowiedni czas życia
źródła. Nie należy też zakładać, że splitowanie tekstu CSV obsługuje
cudzysłowy i escapowanie pól zgodnie z pełnym formatem CSV.

\paragraph{Uwaga na wersję standardu.}
\texttt{std::views::zip} i \texttt{std::views::enumerate} są
konstrukcjami \textbf{C++23}, podobnie jak \texttt{std::ranges::to}.
Nie są dostępne w czystym C++20. W standardzie C++20 wynik można zebrać do kontenera np.
algorytmem \texttt{ranges::copy} z iteratorem wstawiającym
\texttt{std::back\_inserter}.

\section{Materializacja wyniku i koszt przetwarzania}
\label{sec:cpp8-materialize}

Jeśli wynik ma być niezależny od źródła i przechowywany na później,
możemy zamienić potok w kontener. W C++20 nie ma standardowego
\texttt{std::ranges::to<std::vector>()}.

\begin{lstlisting}
#include <algorithm>
#include <iterator>
#include <vector>

std::vector<int> data{1,2,3,4,5,6};
auto even = data | std::views::filter([](int x) {
    return x % 2 == 0;
});

std::vector<int> result;
std::ranges::copy(even, std::back_inserter(result));
// result == {2, 4, 6}
\end{lstlisting}

Taka materializacja wymaga pamięci proporcjonalnej do wielkości wyniku.
Zaletą potoku jest unikanie pośrednich alokacji i często możliwość
\emph{łączenia} transformacji w jednym przebiegu. Nie jest to jednak
gwarancja, że widoki zawsze będą szybsze: koszty iteracji, predykatów,
przepływu sterowania i optymalizacji zależą od konkretnego kodu i
kompilatora. W zadaniach krytycznych wykonuj pomiary.

\section{Własność danych, czas życia i \texttt{borrowed\_range}}
\label{sec:cpp8-lifetime}

To najważniejsza część praktycznego korzystania z Ranges. Widoki
\textbf{nie zawsze posiadają} bazowe dane. W przypadku
\texttt{std::views::all} przeprowadzamy istotne rozróżnienie:

\begin{lstlisting}
std::vector<int> source{2, 4, 6};
auto borrowed = std::views::all(source);
// Refers to source: source must stay alive.

auto owned = std::views::all(std::vector<int>{3, 5, 7});
// In C++20, this can be an owning_view.
// The temporary vector is moved into the view.
\end{lstlisting}

Nie oznacza to, że \emph{każdy} widok utworzony z wartości tymczasowej
przejmuje własność, ani że wszystkie widoki zachowują się tak samo.
Różne adaptory przechowują różne rodzaje obiektów bazowych.

\subsection{Wskaźniki wiszące wewnątrz potoku}

\begin{lstlisting}
auto make_bad_view() {
    std::vector<int> local{1, 2, 3};
    return local | std::views::filter([](int x) {
        return x > 1;
    });
}
// Using the returned view would access destroyed local.
\end{lstlisting}

Ponieważ \texttt{local} jest \emph{lvalue}, adapter przechowuje
nieposiadający danych widok do lokalnego kontenera. Po wyjściu z
funkcji \texttt{local} zostaje zniszczony, a zwrócony widok zawiera
nieaktualne odwołanie. Bezpieczne wyjścia to utworzenie widoku z
przekazanego na własność kontenera (np. obiektu tymczasowego lub
\texttt{std::move(local)}) albo zwrócenie zwykłego kontenera wynikowego.

\begin{lstlisting}
auto make_owning_view() {
    std::vector<int> local{1, 2, 3};
    return std::views::all(std::move(local))
         | std::views::filter([](int x) { return x > 1; });
}
// The returned view owns its vector through owning_view.
\end{lstlisting}

\subsection{Co naprawdę mówi \texttt{borrowed\_range}?}

Koncepcja \texttt{borrowed\_range<R>} oznacza w uproszczeniu, że
iteratory otrzymane z zakresu typu \texttt{R} mogą pozostać ważne
po zniszczeniu \emph{samego obiektu zakresu} (o ile dane, do których
odsyłają, nadal istnieją). Nie jest to synonim ``widok nie posiada
danych'' ani ``wszystkie wskaźniki są zawsze bezpieczne''.

\begin{lstlisting}
#include <concepts>
#include <ranges>
#include <vector>

static_assert(std::ranges::borrowed_range<
    std::ranges::ref_view<std::vector<int>>>);
static_assert(!std::ranges::borrowed_range<
    std::vector<int>>);
\end{lstlisting}

\subsection{\texttt{std::ranges::dangling} jako zabezpieczenie}

Algorytm zwracający iterator do \emph{tymczasowego} kontenera nie może
bezpiecznie zwrócić rzeczywistego iteratora do zniszczonych danych.
Dlatego pewne algorytmy Ranges zwracają specjalny typ
\texttt{std::ranges::dangling}.

\begin{lstlisting}
#include <type_traits>

using Result = decltype(std::ranges::find(
    std::vector<int>{1, 2, 3}, 2));
static_assert(std::same_as<Result, std::ranges::dangling>);
\end{lstlisting}

Nie należy jednak oczekiwać, że mechanizm ten wykryje wszystkie
problemy z czasem życia. Jeżeli zapisujemy nieposiadający danych
widok do lokalnej zmiennej, kompilator może nie zauważyć, że docelowy
kontener umrze wcześniej niż widok.

\section{\texttt{std::span}, \texttt{std::string\_view} i \texttt{subrange}}
\label{sec:cpp8-span}

\texttt{std::span<T>} opisuje ciągły fragment pamięci jako wskaźnik
i rozmiar (lub rozmiar znany w czasie kompilacji). Jest to lekka,
\textbf{nieposiadająca danych} alternatywa dla pary
\texttt{(T*, size)}.

\begin{lstlisting}
#include <span>

void normalize(std::span<double> xs) {
    for (double& x : xs) x /= 100.0;
}

std::vector<double> v{100, 200, 300};
normalize(v);  // v becomes {1, 2, 3}
\end{lstlisting}

\texttt{std::span<const T>} umożliwia odczyt bez modyfikowania
elementów, a \texttt{std::span<T>} pozwala modyfikować elementy, mimo
że sam widok nie posiada pamięci.

\begin{lstlisting}
int a[]{1, 2, 3, 4};
std::span<int> all{a};
auto head = all.first(2);
head[1] = 20; // a[1] == 20
\end{lstlisting}

\texttt{std::string\_view} spełnia podobną rolę dla tekstu; nie
alokuje pamięci na znaki (oraz nie musi kończyć się bajtem zero).
Trzeba uważać, aby nie zapisać go po zniszczeniu źródłowego
\texttt{std::string}.

\begin{lstlisting}
#include <string_view>

std::string text = "abcdef";
std::string_view view{text};
std::string_view prefix = view.substr(0, 3); // "abc"
// Do not use view after text is destroyed or invalidated.
\end{lstlisting}

\texttt{std::ranges::subrange<I, S>} grupuje iterator \texttt{I}
i sentinel \texttt{S} w zakres bez kopiowania elementów. Także tutaj
obowiązuje odpowiedzialność za czas życia wskazywanych danych.

\begin{lstlisting}
auto sub = std::ranges::subrange(v.begin(), v.end());
// sub refers to elements of v; it does not own v.
\end{lstlisting}

\paragraph{Ostrożnie z realokacją.}
\texttt{std::span} do elementów wektora, \texttt{string\_view} do
znaków łańcucha oraz \texttt{subrange} z iteratorów do kontenera mogą
stać się niepoprawne po operacji zmieniającej miejsce przechowywania
danych. Przenoszenie lub realokacja nie są automatycznie śledzone.

\section{Własny widok za pomocą \texttt{view\_interface}}
\label{sec:cpp8-custom-view}

Na końcu zbudujemy \emph{własny widok} generujący kwadraty liczb
całkowitych od zera do \(n-1\). Użyjemy
\texttt{std::ranges::view\_interface}, które dostarcza wygodnych
operacji opartych na \texttt{begin()} i \texttt{end()}. W klasie
pozostają wyłącznie parametry generatora, więc nie ma alokacji pamięci
na elementy sekwencji.

\begin{lstlisting}
#include <concepts>
#include <cstddef>
#include <iterator>
#include <ranges>

struct SquaresView : std::ranges::view_interface<SquaresView> {
    struct iterator {
        using value_type = int;
        using difference_type = std::ptrdiff_t;
        using iterator_concept = std::forward_iterator_tag;

        int index = 0;
        int operator*() const { return index * index; }
        iterator& operator++() {
            ++index;
            return *this;
        }
        iterator operator++(int) {
            auto old = *this;
            ++*this;
            return old;
        }
        bool operator==(const iterator&) const = default;
    };

    int n = 0;
    SquaresView() = default;
    explicit SquaresView(int count) : n(count) {}

    iterator begin() const { return {0}; }
    iterator end() const { return {n}; }
};

static_assert(std::ranges::view<SquaresView>);
static_assert(std::ranges::forward_range<SquaresView>);
\end{lstlisting}

\begin{lstlisting}
SquaresView squares{6};
for (int x : squares) { /* 0, 1, 4, 9, 16, 25 */ }

// Our view works with standard adapters:
auto even_squares = squares | std::views::filter(
    [](int x) { return x % 2 == 0; });
// Iteration yields 0, 4, 16.
\end{lstlisting}

Konstruktor z argumentem \texttt{count} powinien otrzymać wartość
nieujemną. Ponieważ liczymy \texttt{index * index} w typie
\texttt{int}, należy ponadto ograniczyć rozmiar zakresu, aby uniknąć
przepełnienia. W rozwiązaniu produkcyjnym należałoby te warunki
weryfikować oraz dopracować obsługę szerszych typów całkowitych.

\paragraph{Widok a adapter.}
\texttt{SquaresView} jest własnym \emph{widokiem}, który można
przekazać do standardowych adapterów. Nie jest jeszcze własnym
\emph{range adaptor closure}, czyli obiektem pozwalającym zapisać
\texttt{some\_range | squares\_adaptor}. Implementacja takiego
adaptera stanowi osobny poziom abstrakcji i nie jest konieczna do
budowania własnych widoków.

\section{Różnica między \texttt{transform} i \texttt{views::transform}}

Poniższe operacje mogą prowadzić do takich samych liczb, ale różnią
się \emph{czasem wykonania} i \emph{sposobem przechowywania wyniku}:

\begin{lstlisting}
std::vector<int> source{1, 2, 3};
std::vector<int> eager;
std::ranges::transform(source,
    std::back_inserter(eager),
    [](int x) { return x * x; });
// eager owns {1, 4, 9} right now

auto lazy = source | std::views::transform(
    [](int x) { return x * x; });
// lazy does not own a vector of squared values
\end{lstlisting}

Wariant pierwszy wymaga pamięci na trzy elementy i kończy obliczenia
przed kolejną instrukcją. Wariant drugi może obliczać wartości dopiero
przy dereferencji iteratorów. Po zmianie \texttt{source} kolejne
odczyty widoku \texttt{lazy} mogą odzwierciedlać nowe dane.

\section{Projektowanie API z wykorzystaniem Ranges}
\label{sec:cpp8-api}

Jeżeli funkcja tylko przegląda elementy dowolnego źródła, preferuj
parametr ograniczony koncepcją \texttt{input\_range}. Dla wielokrotnych
przejść korzystaj z \texttt{forward\_range}. Jeśli potrzebujesz
ciągłego bloku pamięci i znasz typ elementów, wygodnym wyborem jest
\texttt{std::span<T>}.

\begin{lstlisting}
template <std::ranges::input_range R>
void print_each(R&& r) {
    for (auto&& x : r) {
        // Print or consume x.
    }
}

void process_buffer(std::span<const double> values) {
    for (double value : values) {
        // Read numeric values.
    }
}
\end{lstlisting}

Zależnie od celu można zwracać widok albo kontener. Widok jest
wygodny, gdy chcemy dalej budować potok i mamy kontrolę nad czasem
życia źródła. Kontener jest często bezpieczniejszy, gdy zwracany
wynik ma istnieć niezależnie od wszystkich argumentów.

\section{Najczęstsze błędy i praktyczne reguły}
\label{sec:cpp8-pitfalls}

\begin{enumerate}
\item \textbf{Dereferencja \texttt{end()}.} Iterator końcowy jest
granicą, a nie elementem. Przed \texttt{*it} trzeba mieć
\texttt{it != end}.

\item \textbf{Mylenie kategorii iteratorów.} Lista nie obsługuje
skoków w stałym czasie. Nie można użyć dla niej algorytmu
\texttt{ranges::sort}; należy wywołać metodę \texttt{list::sort}.

\item \textbf{Traktowanie widoku jak kontenera.} \texttt{views::filter}
nie tworzy kopii wybranych elementów. Zmiany i zniszczenie źródła
mogą wpływać na wynik.

\item \textbf{Wiszące referencje w lambdach.} Zwrócony widok
z predykatem przechwytującym lokalny obiekt przez referencję może
być niebezpieczny nawet wtedy, gdy sam kontener ma odpowiednio długi
czas życia.

\item \textbf{Efekty uboczne w predykacie.} Liczba i kolejność
wywołań w leniwym potoku nie są odpowiednią podstawą działania
programu; predykaty filtrowania powinny zachować stabilną semantykę.

\item \textbf{Zakładanie stałego kosztu operacji.} Obliczenie
\texttt{distance} lub pierwszego pasującego elementu może wymagać
przejrzenia dużej części źródła.

\item \textbf{Mieszanie standardów.} \texttt{ranges::to},
\texttt{views::zip} i \texttt{views::enumerate} nie należą do C++20.

\item \textbf{Używanie nieważnych iteratorów.} Realokacja wektora
lub usuwanie elementu może unieważnić iteratory także wtedy, gdy
znajdują się wewnątrz wcześniej utworzonego widoku.
\end{enumerate}

\section{Której konstrukcji użyć?}

\begin{center}
\small
\begin{tabular}{|p{6.0cm}|p{6.5cm}|}
\hline
\textbf{Potrzeba} & \textbf{Zalecany wybór} \\
\hline
Przejście po zwykłym kontenerze & Pętla \texttt{for (auto\& x : v)} \\
\hline
Jednorazowe sortowanie albo wyszukiwanie & \texttt{std::ranges::sort}, \texttt{std::ranges::find} \\
\hline
Leniwe filtrowanie i przekształcanie & \texttt{std::views::filter}, \texttt{std::views::transform} \\
\hline
Generowanie sekwencji indeksów & \texttt{std::views::iota} \\
\hline
Odczyt lub modyfikacja ciągłego bufora & \texttt{std::span<T>} (bez własności) \\
\hline
Lekki odczyt fragmentu tekstu & \texttt{std::string\_view} (bez własności) \\
\hline
Zakres wyznaczony iteratorami & \texttt{std::ranges::subrange} \\
\hline
Wynik niezależny od źródła & Kontener, np. \texttt{std::vector} \\
\hline
Własny generator wartości & Typ modelujący \texttt{std::ranges::view} \\
\hline
Wydajny ogólny algorytm generyczny & Szablon z \texttt{std::ranges::input\_range} lub silniejszą koncepcją \\
\hline
\end{tabular}
\end{center}

\paragraph{Pytania kontrolne.}
\begin{enumerate}
\item Dlaczego iterator \texttt{end()} nie jest elementem kolekcji?
\item Co odróżnia \texttt{forward\_iterator} od \texttt{input\_iterator}?
\item Czy \texttt{std::ranges::sort} działa dla \texttt{std::list}? Dlaczego?
\item Po co \texttt{std::ranges} dopuszcza sentinel innego typu niż iterator?
\item W którym momencie \texttt{views::transform} wykonuje funkcję?
\item Dlaczego \texttt{views::filter} nie powinien zmieniać elementów, od których zależy jego predykat?
\item Czym różni się \texttt{views::all(lvalue)} od \texttt{views::all(rvalue)} dla \texttt{std::vector}?
\item Co gwarantuje, a czego nie gwarantuje \texttt{borrowed\_range}?
\item Dlaczego \texttt{ranges::find(std::vector\{\dots\}, x)} może zwrócić \texttt{dangling}?
\item Jak zebrać elementy \texttt{filter\_view} do wektora w czystym C++20?
\item Czym różni się \texttt{std::span<const T>} od \texttt{std::vector<T>}?
\item Co należy zaimplementować, aby własny obiekt był \texttt{forward\_range}?
\end{enumerate}

\paragraph{Dalsza lektura.}
Dobrym punktem odniesienia jest dokumentacja \emph{cppreference}:
\emph{Ranges library}, \emph{Iterator library}, \emph{Range adaptors},
\emph{std::ranges::borrowed\_range}, \emph{std::ranges::dangling}
oraz \emph{std::ranges::view\_interface}. Warto testować kod z opcjami kompilatora
\texttt{-std=c++20} i \texttt{-Wall -Wextra -Wpedantic},
a także korzystać z AddressSanitizer przy badaniu czasu życia danych.


\chapter[Iteratory, Ranges i Views]{Iteratory, zakresy i widoki: \textit{Ranges} oraz \textit{Views} w C++20}
\label{chap:ranges-views-cpp}

\section{Dlaczego wprowadzono Ranges?}

Klasyczne algorytmy STL operują na \emph{parze iteratorów}: pierwszym
wskazującym początek danych i drugim wyznaczającym koniec. Dzięki temu
\texttt{std::sort} lub \texttt{std::find} mogą współpracować z wieloma
kontenerami. Wymaga to jednak przekazywania obu granic i dodatkowej
ostrożności, aby pochodziły z tego samego zakresu.

\begin{lstlisting}
#include <algorithm>
#include <vector>

std::vector<int> xs{4, 1, 3, 2};
std::sort(xs.begin(), xs.end()); // STL before C++20
std::ranges::sort(xs);          // C++20
\end{lstlisting}

Nowy model pozwala przekazać \emph{zakres} (\textit{range}) jako jeden
obiekt oraz formułować wymagania wobec niego poprzez \emph{koncepcje}
(\textit{concepts}, rozdział~5). Oddzielną innowacją są \emph{widoki}
(\textit{views}): lekkie reprezentacje danych, często przekształcane
\textbf{leniwe}, to znaczy dopiero podczas ich przeglądania.

\begin{lstlisting}
#include <ranges>

std::vector<int> xs{1, 2, 3, 4, 5, 6};
auto pipeline = xs
    | std::views::filter([](int x) { return x % 2 == 0; })
    | std::views::transform([](int x) { return x * x; });

// No vector of squares has been created.
// Iteration produces 4, 16, 36.
for (int x : pipeline) { /* use x */ }
\end{lstlisting}

\paragraph{Założenia.}
W całym rozdziale przyjmujemy standard \textbf{C++20} oraz nagłówki
\texttt{<algorithm>}, \texttt{<iterator>}, \texttt{<ranges>},
\texttt{<concepts>}, \texttt{<vector>} i pozostałe nagłówki wymienione
w przykładach. W dalszych sekcjach posłużymy się semantyką przenoszenia
(rozdział~6) i lambdami (rozdział~7).

\section{Iteratory: uogólnienie wskaźników}
\label{sec:cpp8-iterators}

\subsection{Operacje \texttt{begin}, \texttt{end}, \texttt{*} i \texttt{++}}

Iterator jest obiektem umożliwiającym poruszanie się po elementach
zakresu. Wskaźnik \texttt{int*} jest przykładem iteratora, ale iterator
może także przechowywać indeks, wskaźnik do węzła listy lub stan
wewnętrzny generatora. Podstawowe operacje przedstawia program:

\begin{lstlisting}
#include <vector>

std::vector<int> v{10, 20, 30};
auto it = v.begin();     // points to 10
int a = *it;            // 10
++it;                   // points to 20
int b = *it;            // 20
bool ok = it != v.end();
\end{lstlisting}

Iterator \texttt{end()} zwykle oznacza pozycję \emph{za ostatnim}
elementem; nie wolno go dereferencjonować. Dzięki temu pusty zakres
spełnia \texttt{begin()==end()}, a przechodzenie po elementach ma
jednolity schemat:

\begin{lstlisting}
for (auto it = v.begin(); it != v.end(); ++it) {
    // *it is a valid element
}
\end{lstlisting}

\paragraph{Pętla zakresowa.}
Zapis \texttt{for (auto\& x : v)} jest wygodniejszą postacią iteracji,
ale w dalszym ciągu wykorzystuje operacje początku, końca,
porównania, zwiększania i dereferencji. Aby \emph{modyfikować} elementy,
stosujemy \texttt{auto\&}; aby uniknąć kopiowania bez modyfikowania,
\texttt{const auto\&}; dla niewielkich wartości liczbowych często
wystarczy \texttt{auto}.

\subsection{Iteratory stałe i nieważność iteratorów}

\begin{lstlisting}
std::vector<int> v{1, 2, 3};
auto m = v.begin();   // can modify *m
auto c = v.cbegin();  // cannot modify *c
*m = 10;
// *c = 20;         // compilation error
\end{lstlisting}

\texttt{const\_iterator} chroni elementy przed modyfikacją poprzez
ten iterator; nie oznacza, że cały kontener jest niezmienny. Równie
ważna jest \emph{nieważność iteratorów} (\textit{iterator invalidation}).
Na przykład ponowna alokacja pamięci przez \texttt{std::vector::push\_back}
może unieważnić \textbf{wszystkie} iteratory, wskaźniki i referencje do
jego elementów.

\begin{lstlisting}
std::vector<int> v{1, 2, 3};
auto it = v.begin();
v.push_back(4);       // may reallocate storage
// *it;              // unsafe if reallocation happened
it = v.begin();      // obtain a fresh iterator
\end{lstlisting}

Reguły zależą od kontenera i operacji. W \texttt{std::list} wstawienie
elementu nie unieważnia iteratorów do pozostałych elementów, ale nie
oznacza to, że można dalej używać iteratora do elementu usuniętego.

\section{Kategorie iteratorów i ich wymagania}
\label{sec:cpp8-iterator-categories}

Zestaw koncepcji C++20 opisuje operacje i własności iteratorów, a nie
jedynie nazwy klas. Najważniejsze kategorie zestawia tabela.

\begin{center}
\small
\begin{tabular}{|p{5.2cm}|p{7.3cm}|}
\hline
\textbf{Koncepcja} & \textbf{Możliwości i przykłady} \\
\hline
\texttt{input\_iterator} & Odczytywanie kolejnych wartości; przeglądanie może być jednokrotne (np. strumień wejściowy). \\
\hline
\texttt{output\_iterator} & Zapisywanie kolejnych wartości do celu (np. \texttt{back\_insert\_iterator}). \\
\hline
\texttt{forward\_iterator} & Wielokrotne przejścia, porównywanie iteratorów, krok do przodu. \\
\hline
\texttt{bidirectional\_iterator} & Dodatkowo krok wstecz, czyli \texttt{--it}; np. \texttt{list}. \\
\hline
\texttt{random\_access\_iterator} & Indeksowanie i skoki \texttt{it+n} w stałym czasie; np. \texttt{deque}. \\
\hline
\texttt{contiguous\_iterator} & Elementy leżą kolejno w pamięci; np. iterator \texttt{vector<int>} i surowy wskaźnik. \\
\hline
\end{tabular}
\end{center}

Hierarchia wymagań dla iteratorów odczytujących jest w przybliżeniu
następująca (każda kolejna koncepcja jest silniejsza):

\begin{center}
\texttt{input $\Rightarrow$ forward $\Rightarrow$ bidirectional}
\\[2pt]
\texttt{$\Rightarrow$ random\_access $\Rightarrow$ contiguous}.
\end{center}

Koncepcja \texttt{output\_iterator} dotyczy osobnej rodziny wymagań.
Nowe koncepcje nie są w każdym szczególe identyczne ze starymi
kategoriami \emph{LegacyIterator}. Na przykład model
\texttt{forward\_iterator} nie wymaga zwracania referencji
przy dereferencji.

\begin{lstlisting}
#include <concepts>
#include <iterator>
#include <list>
#include <vector>

static_assert(std::contiguous_iterator<int*>);
static_assert(std::random_access_iterator<
    std::vector<int>::iterator>);
static_assert(std::bidirectional_iterator<
    std::list<int>::iterator>);
static_assert(!std::random_access_iterator<
    std::list<int>::iterator>);
\end{lstlisting}

\paragraph{Dlaczego to ważne?}
\texttt{std::ranges::sort} wymaga iteratorów o dostępie swobodnym,
więc można użyć go do \texttt{std::vector}, ale \textbf{nie} do
\texttt{std::list}. Listę sortuje się metodą \texttt{list::sort()}.
Wybór algorytmu jest zatem ograniczony kategorią iteratora.

\section{Własny iterator i własny zakres}
\label{sec:cpp8-own-iterator}

Zbudujemy iterator generujący kolejne liczby całkowite. Nie będzie on
przechowywać ani wskaźnika, ani kontenera: całą potrzebną informacją
jest bieżąca wartość. Zdefiniujemy \emph{iterator do przodu}
(\texttt{forward\_iterator}).

\begin{lstlisting}
#include <concepts>
#include <cstddef>
#include <iterator>
#include <ranges>

struct CountingIterator {
    using value_type = int;
    using difference_type = std::ptrdiff_t;
    using iterator_concept = std::forward_iterator_tag;

    int current = 0;

    int operator*() const { return current; }

    CountingIterator& operator++() {
        ++current;
        return *this;
    }
    CountingIterator operator++(int) {
        auto old = *this;
        ++*this;
        return old;
    }
    bool operator==(const CountingIterator&) const = default;
};

struct CountingRange {
    int first = 0;
    int last = 0;

    CountingIterator begin() const { return {first}; }
    CountingIterator end() const { return {last}; }
};

static_assert(std::forward_iterator<CountingIterator>);
static_assert(std::ranges::forward_range<CountingRange>);
\end{lstlisting}

Teraz kod \texttt{for (int i : CountingRange\{3, 7\})} wypisze
\texttt{3, 4, 5, 6}. Funkcje \texttt{begin()} i \texttt{end()}
wyznaczają \emph{półotwarty} zakres \([3,7)\). Prawidłowe użycie
zakłada \texttt{first <= last}; dla \texttt{first > last} nasz
minimalny iterator nie zbiegnie do wartości końcowej.

\paragraph{Co trzeba zagwarantować?}
Samo przejście \texttt{static\_assert} nie dowodzi wszystkich
semantycznych własności iteratora. Dla \texttt{forward\_iterator}
ważna jest m.in. własność \emph{multi-pass}: kopiowanie iteratora
i przejście jedną kopią nie powinno zmieniać wartości zwracanych
przez drugą. Nasz iterator ją spełnia, ponieważ każda kopia ma własne
pole \texttt{current}.

\section{Zwykłe \textit{range} i \textit{sentinel}}
\label{sec:cpp8-range-sentinel}

\emph{Zakres} jest obiektem, dla którego można otrzymać początek
\texttt{std::ranges::begin(r)} i koniec \texttt{std::ranges::end(r)}.
Typ wyniku obu wywołań \textbf{nie musi być identyczny}. Obiekt
wyznaczający koniec to \emph{sentinel}; musi dać się porównać z
iteratorem, ale sam nie musi obsługiwać jego pozostałych operacji.

\begin{lstlisting}
#include <ranges>
#include <vector>

std::vector<int> xs{1, 2, 3};
auto first = std::ranges::begin(xs);
auto last = std::ranges::end(xs);
while (first != last) { ++first; }
\end{lstlisting}

W tym przypadku obie granice mają ten sam typ: zakres spełnia
koncepcję \texttt{common\_range}. W innych zakresach koniec może
wyznaczać osobny obiekt sentinela. Odpowiednia koncepcja
\texttt{std::sentinel\_for} określa warunki porównywania granic:

\begin{lstlisting}
static_assert(std::ranges::common_range<decltype(xs)>);
static_assert(std::sentinel_for<
    decltype(last), decltype(first)>);
\end{lstlisting}

Przykładem osobnego typu końca jest \texttt{std::default\_sentinel\_t}.

\paragraph{Rozmiar a liczba kroków.}
Operacja \texttt{std::ranges::size(r)} jest dostępna dla zakresów
spełniających \texttt{sized\_range}. Nie każdy zakres ma tanią operację
wyznaczenia rozmiaru. Funkcja \texttt{distance(r)} z przestrzeni
\texttt{std::ranges} może wymagać
\emph{przejścia przez wszystkie elementy}. Dla niektórych
jednokrotnych źródeł danych przejście oznacza również ich zużycie.

\section{Koncepcje zakresów w praktyce}

Własności zakresów sprawdzamy za pomocą koncepcji z rozdziału~5.
Najważniejsze z nich to:
\begin{itemize}
\item \texttt{input\_range} i \texttt{forward\_range},
\item \texttt{bidirectional\_range} i \texttt{random\_access\_range},
\item \texttt{contiguous\_range}, \texttt{common\_range} i \texttt{sized\_range}.
\end{itemize}

\begin{lstlisting}
#include <ranges>
#include <vector>

template <std::ranges::input_range R>
long long sum_integers(R&& range) {
    long long result = 0;
    for (auto value : range) result += value;
    return result;
}

std::vector<int> v{1, 2, 3};
long long s = sum_integers(v); // 6
\end{lstlisting}

Parametr \texttt{R\&\&} jest tutaj \emph{referencją przekazującą}
(\textit{forwarding reference}), a nie żądaniem przejęcia własności.
Funkcja obsługuje także źródła jednokrotne, gdyż przechodzi zakres tylko
raz. Dla naprawdę uniwersalnej funkcji trzeba osobno wyrazić warunki
na typ elementu i wykonywaną na nim operację dodawania.

\section{Algorytmy \texttt{std::ranges} i projekcje}
\label{sec:cpp8-algorithms}

Funkcje z rodziny \texttt{std::ranges}, np. \texttt{sort},
\texttt{find}, \texttt{copy} i \texttt{transform}, przyjmują
bezpośrednio zakres albo parę iterator--sentinel.
Ich wykonanie jest na ogół \emph{natychmiastowe} (w odróżnieniu od
samych widoków \texttt{views::transform}).

\begin{lstlisting}
#include <algorithm>
#include <iterator>
#include <vector>

std::vector<int> v{4, 2, 3, 1};
std::ranges::sort(v);  // modifies v immediately

auto pos = std::ranges::find(v, 3);
if (pos != v.end()) { *pos = 30; }

std::vector<int> out;
std::ranges::copy(v, std::back_inserter(out));
\end{lstlisting}

\subsection{Projekcja: sortowanie po wskazanym polu}

Algorytmy Ranges często posiadają argument \emph{projekcji}: funkcję
wywoływaną na elemencie przed porównaniem albo wyszukiwaniem. Dzięki
temu zamiast pisać własny komparator, można wskazać pole struktury.

\begin{lstlisting}
#include <string>
#include <vector>

struct Person {
    std::string name;
    int age;
};

std::vector<Person> people{
    {"Ada", 40}, {"Ewa", 25}, {"Jan", 32}
};

std::ranges::sort(people, {}, &Person::age);
// ascending ages: 25, 32, 40

auto it = std::ranges::find(people, 32, &Person::age);
// it->name == "Jan"
\end{lstlisting}

W wywołaniu \texttt{sort(people, \{\}, \&Person::age)} puste
nawiasy oznaczają domyślny komparator \texttt{std::ranges::less},
a wskaźnik do składowej to projekcja. Standard używa mechanizmu
\texttt{std::invoke}, omówionego w rozdziale~7.

\section{Widoki: \textit{view} nie jest synonimem kontenera}
\label{sec:cpp8-views}

\emph{Widok} jest zakresem przystosowanym do taniego tworzenia i
kopiowania lub przenoszenia \emph{samego obiektu widoku}. Nie musi mieć
własnych danych: może przechowywać referencję do kontenera, mały zestaw
iteratorów albo parametry generowania wartości.

\begin{lstlisting}
#include <ranges>
#include <vector>

std::vector<int> xs{10, 20, 30, 40};
auto first_two = xs | std::views::take(2);
// first_two refers to the original xs
\end{lstlisting}

Stworzenie widoku nie jest równoznaczne ze skopiowaniem elementów do
nowego kontenera. Odwołanie do danych źródłowych wymaga, by pozostawały
one ważne przez cały czas korzystania z widoku. Szczegółowe reguły
własności omówimy dalej.

\paragraph{Dwa znaczenia nazwy \texttt{views}.}
\texttt{std::views} jest aliasem przestrzeni nazw z obiektami
adaptującymi zakresy. \texttt{std::ranges::view} to natomiast
\emph{koncepcja} spełniana przez dany typ. Nie należy mylić ich
z pojęciem \texttt{std::ranges::range}.

\section{Leniwe potoki \texttt{filter}, \texttt{transform}, \texttt{take}}
\label{sec:cpp8-pipelines}

Najczęściej stosowane widoki łączymy za pomocą operatora \texttt{|}.
Adaptory przekazują sobie zakresy bez potrzeby budowania pośrednich
wektorów:

\begin{lstlisting}
std::vector<int> xs{1,2,3,4,5,6,7,8,9,10};

auto pipeline = xs
    | std::views::filter([](int x) {
          return x % 2 == 0;
      })
    | std::views::transform([](int x) {
          return x * x;
      })
    | std::views::take(3);

// Iteration produces 4, 16, 36.
\end{lstlisting}

\texttt{filter} pomija elementy niespełniające predykatu,
\texttt{transform} wyznacza nową wartość przy dereferencji,
a \texttt{take(3)} ogranicza liczbę wyników. Algorytm
\texttt{std::ranges::transform} działa od razu; natomiast
utworzenie widoku \texttt{views::transform} nie przelicza
wszystkich wyników.

\subsection{Kiedy następują obliczenia?}

\begin{lstlisting}
int calls = 0;
auto squares = xs | std::views::transform([&](int x) {
    ++calls;
    return x * x;
});
// calls == 0 here
int first = *squares.begin();
// calls has increased (transform was evaluated)
\end{lstlisting}

W tym przykładzie pokazujemy moment obliczenia, a nie zalecany styl
projektowania: \textbf{predykaty i transformacje powinny zasadniczo
unikać efektów ubocznych}. Ponowne przejście widoku, ponowna
dereferencja lub implementacyjnie dopuszczone buforowanie może
powodować inne liczby wywołań, niż sugerowałaby pojedyncza pętla.
Nie należy projektować logiki programu na podstawie rzekomej liczby
wywołań funkcji przez widok.

\paragraph{Ważny szczegół \texttt{filter\_view}.}
Dla pewnych kategorii zakresu \texttt{filter\_view} może zapamiętywać
pierwszy znaleziony element. Ponadto modyfikowanie elementu tak,
aby przestał spełniać predykat filtrowania, narusza wymagania
\texttt{filter\_view} i może prowadzić do niezdefiniowanego
zachowania. Predykat powinien być stabilny wobec przebiegu iteracji.

\section{Pozostałe przydatne adaptory C++20}
\label{sec:cpp8-adaptors}

\subsection{\texttt{iota}: wartości wyznaczane bez tablicy}

\begin{lstlisting}
auto ids = std::views::iota(1, 5);
// Iteration produces 1, 2, 3, 4.

// This is an unbounded range, limited by take(5):
auto first_five = std::views::iota(1)
                | std::views::take(5);
// Produces 1, 2, 3, 4, 5.
\end{lstlisting}

\texttt{iota} przydaje się do indeksowania lub generowania sekwencji.
Dla nieograniczonego widoku trzeba pamiętać o warunku zakończenia;
ponadto operacje arytmetyczne nie mogą przepełniać typu całkowitego
ze znakiem.

\subsection{\texttt{drop}, \texttt{take} i \texttt{reverse}}

\begin{lstlisting}
std::vector<int> data{1, 2, 3, 4, 5};
auto middle = data | std::views::drop(1)
                   | std::views::take(3);
// 2, 3, 4

auto backwards = data | std::views::reverse;
// 5, 4, 3, 2, 1
\end{lstlisting}

\texttt{reverse} wymaga co najmniej zakresu dwukierunkowego;
\texttt{take} i \texttt{drop} są przydatne także dla sekwencji
generowanych leniwie. Dla niektórych widoków koszt \texttt{begin()}
lub dojścia do pierwszego wyniku może być większy niż stały.

\subsection{\texttt{keys}, \texttt{values} i \texttt{split}}

\begin{lstlisting}
std::vector<std::pair<int, std::string>> pairs{
    {1, "one"}, {2, "two"}
};
for (int key : pairs | std::views::keys) { /* 1, 2 */ }
for (const auto& value : pairs | std::views::values) {
    /* "one", "two" */
}

std::string csv = "aa,b,cc";
for (auto part : csv | std::views::split(',')) {
    std::string piece(part.begin(), part.end());
    // Copy each part into an owning string if needed.
}
\end{lstlisting}

\texttt{split} nie tworzy automatycznie \texttt{std::string} dla
każdego fragmentu; otrzymujemy zakresy wskazujące fragmenty źródła.
Zachowując wynik na później, trzeba zapewnić odpowiedni czas życia
źródła. Nie należy też zakładać, że splitowanie tekstu CSV obsługuje
cudzysłowy i escapowanie pól zgodnie z pełnym formatem CSV.

\paragraph{Uwaga na wersję standardu.}
\texttt{std::views::zip} i \texttt{std::views::enumerate} są
konstrukcjami \textbf{C++23}, podobnie jak \texttt{std::ranges::to}.
Nie są dostępne w czystym C++20. W standardzie C++20 wynik można zebrać do kontenera np.
algorytmem \texttt{ranges::copy} z iteratorem wstawiającym
\texttt{std::back\_inserter}.

\section{Materializacja wyniku i koszt przetwarzania}
\label{sec:cpp8-materialize}

Jeśli wynik ma być niezależny od źródła i przechowywany na później,
możemy zamienić potok w kontener. W C++20 nie ma standardowego
\texttt{std::ranges::to<std::vector>()}.

\begin{lstlisting}
#include <algorithm>
#include <iterator>
#include <vector>

std::vector<int> data{1,2,3,4,5,6};
auto even = data | std::views::filter([](int x) {
    return x % 2 == 0;
});

std::vector<int> result;
std::ranges::copy(even, std::back_inserter(result));
// result == {2, 4, 6}
\end{lstlisting}

Taka materializacja wymaga pamięci proporcjonalnej do wielkości wyniku.
Zaletą potoku jest unikanie pośrednich alokacji i często możliwość
\emph{łączenia} transformacji w jednym przebiegu. Nie jest to jednak
gwarancja, że widoki zawsze będą szybsze: koszty iteracji, predykatów,
przepływu sterowania i optymalizacji zależą od konkretnego kodu i
kompilatora. W zadaniach krytycznych wykonuj pomiary.

\section{Własność danych, czas życia i \texttt{borrowed\_range}}
\label{sec:cpp8-lifetime}

To najważniejsza część praktycznego korzystania z Ranges. Widoki
\textbf{nie zawsze posiadają} bazowe dane. W przypadku
\texttt{std::views::all} przeprowadzamy istotne rozróżnienie:

\begin{lstlisting}
std::vector<int> source{2, 4, 6};
auto borrowed = std::views::all(source);
// Refers to source: source must stay alive.

auto owned = std::views::all(std::vector<int>{3, 5, 7});
// In C++20, this can be an owning_view.
// The temporary vector is moved into the view.
\end{lstlisting}

Nie oznacza to, że \emph{każdy} widok utworzony z wartości tymczasowej
przejmuje własność, ani że wszystkie widoki zachowują się tak samo.
Różne adaptory przechowują różne rodzaje obiektów bazowych.

\subsection{Wskaźniki wiszące wewnątrz potoku}

\begin{lstlisting}
auto make_bad_view() {
    std::vector<int> local{1, 2, 3};
    return local | std::views::filter([](int x) {
        return x > 1;
    });
}
// Using the returned view would access destroyed local.
\end{lstlisting}

Ponieważ \texttt{local} jest \emph{lvalue}, adapter przechowuje
nieposiadający danych widok do lokalnego kontenera. Po wyjściu z
funkcji \texttt{local} zostaje zniszczony, a zwrócony widok zawiera
nieaktualne odwołanie. Bezpieczne wyjścia to utworzenie widoku z
przekazanego na własność kontenera (np. obiektu tymczasowego lub
\texttt{std::move(local)}) albo zwrócenie zwykłego kontenera wynikowego.

\begin{lstlisting}
auto make_owning_view() {
    std::vector<int> local{1, 2, 3};
    return std::views::all(std::move(local))
         | std::views::filter([](int x) { return x > 1; });
}
// The returned view owns its vector through owning_view.
\end{lstlisting}

\subsection{Co naprawdę mówi \texttt{borrowed\_range}?}

Koncepcja \texttt{borrowed\_range<R>} oznacza w uproszczeniu, że
iteratory otrzymane z zakresu typu \texttt{R} mogą pozostać ważne
po zniszczeniu \emph{samego obiektu zakresu} (o ile dane, do których
odsyłają, nadal istnieją). Nie jest to synonim ``widok nie posiada
danych'' ani ``wszystkie wskaźniki są zawsze bezpieczne''.

\begin{lstlisting}
#include <concepts>
#include <ranges>
#include <vector>

static_assert(std::ranges::borrowed_range<
    std::ranges::ref_view<std::vector<int>>>);
static_assert(!std::ranges::borrowed_range<
    std::vector<int>>);
\end{lstlisting}

\subsection{\texttt{std::ranges::dangling} jako zabezpieczenie}

Algorytm zwracający iterator do \emph{tymczasowego} kontenera nie może
bezpiecznie zwrócić rzeczywistego iteratora do zniszczonych danych.
Dlatego pewne algorytmy Ranges zwracają specjalny typ
\texttt{std::ranges::dangling}.

\begin{lstlisting}
#include <type_traits>

using Result = decltype(std::ranges::find(
    std::vector<int>{1, 2, 3}, 2));
static_assert(std::same_as<Result, std::ranges::dangling>);
\end{lstlisting}

Nie należy jednak oczekiwać, że mechanizm ten wykryje wszystkie
problemy z czasem życia. Jeżeli zapisujemy nieposiadający danych
widok do lokalnej zmiennej, kompilator może nie zauważyć, że docelowy
kontener umrze wcześniej niż widok.

\section{\texttt{std::span}, \texttt{std::string\_view} i \texttt{subrange}}
\label{sec:cpp8-span}

\texttt{std::span<T>} opisuje ciągły fragment pamięci jako wskaźnik
i rozmiar (lub rozmiar znany w czasie kompilacji). Jest to lekka,
\textbf{nieposiadająca danych} alternatywa dla pary
\texttt{(T*, size)}.

\begin{lstlisting}
#include <span>

void normalize(std::span<double> xs) {
    for (double& x : xs) x /= 100.0;
}

std::vector<double> v{100, 200, 300};
normalize(v);  // v becomes {1, 2, 3}
\end{lstlisting}

\texttt{std::span<const T>} umożliwia odczyt bez modyfikowania
elementów, a \texttt{std::span<T>} pozwala modyfikować elementy, mimo
że sam widok nie posiada pamięci.

\begin{lstlisting}
int a[]{1, 2, 3, 4};
std::span<int> all{a};
auto head = all.first(2);
head[1] = 20; // a[1] == 20
\end{lstlisting}

\texttt{std::string\_view} spełnia podobną rolę dla tekstu; nie
alokuje pamięci na znaki (oraz nie musi kończyć się bajtem zero).
Trzeba uważać, aby nie zapisać go po zniszczeniu źródłowego
\texttt{std::string}.

\begin{lstlisting}
#include <string_view>

std::string text = "abcdef";
std::string_view view{text};
std::string_view prefix = view.substr(0, 3); // "abc"
// Do not use view after text is destroyed or invalidated.
\end{lstlisting}

\texttt{std::ranges::subrange<I, S>} grupuje iterator \texttt{I}
i sentinel \texttt{S} w zakres bez kopiowania elementów. Także tutaj
obowiązuje odpowiedzialność za czas życia wskazywanych danych.

\begin{lstlisting}
auto sub = std::ranges::subrange(v.begin(), v.end());
// sub refers to elements of v; it does not own v.
\end{lstlisting}

\paragraph{Ostrożnie z realokacją.}
\texttt{std::span} do elementów wektora, \texttt{string\_view} do
znaków łańcucha oraz \texttt{subrange} z iteratorów do kontenera mogą
stać się niepoprawne po operacji zmieniającej miejsce przechowywania
danych. Przenoszenie lub realokacja nie są automatycznie śledzone.

\section{Własny widok za pomocą \texttt{view\_interface}}
\label{sec:cpp8-custom-view}

Na końcu zbudujemy \emph{własny widok} generujący kwadraty liczb
całkowitych od zera do \(n-1\). Użyjemy
\texttt{std::ranges::view\_interface}, które dostarcza wygodnych
operacji opartych na \texttt{begin()} i \texttt{end()}. W klasie
pozostają wyłącznie parametry generatora, więc nie ma alokacji pamięci
na elementy sekwencji.

\begin{lstlisting}
#include <concepts>
#include <cstddef>
#include <iterator>
#include <ranges>

struct SquaresView : std::ranges::view_interface<SquaresView> {
    struct iterator {
        using value_type = int;
        using difference_type = std::ptrdiff_t;
        using iterator_concept = std::forward_iterator_tag;

        int index = 0;
        int operator*() const { return index * index; }
        iterator& operator++() {
            ++index;
            return *this;
        }
        iterator operator++(int) {
            auto old = *this;
            ++*this;
            return old;
        }
        bool operator==(const iterator&) const = default;
    };

    int n = 0;
    SquaresView() = default;
    explicit SquaresView(int count) : n(count) {}

    iterator begin() const { return {0}; }
    iterator end() const { return {n}; }
};

static_assert(std::ranges::view<SquaresView>);
static_assert(std::ranges::forward_range<SquaresView>);
\end{lstlisting}

\begin{lstlisting}
SquaresView squares{6};
for (int x : squares) { /* 0, 1, 4, 9, 16, 25 */ }

// Our view works with standard adapters:
auto even_squares = squares | std::views::filter(
    [](int x) { return x % 2 == 0; });
// Iteration yields 0, 4, 16.
\end{lstlisting}

Konstruktor z argumentem \texttt{count} powinien otrzymać wartość
nieujemną. Ponieważ liczymy \texttt{index * index} w typie
\texttt{int}, należy ponadto ograniczyć rozmiar zakresu, aby uniknąć
przepełnienia. W rozwiązaniu produkcyjnym należałoby te warunki
weryfikować oraz dopracować obsługę szerszych typów całkowitych.

\paragraph{Widok a adapter.}
\texttt{SquaresView} jest własnym \emph{widokiem}, który można
przekazać do standardowych adapterów. Nie jest jeszcze własnym
\emph{range adaptor closure}, czyli obiektem pozwalającym zapisać
\texttt{some\_range | squares\_adaptor}. Implementacja takiego
adaptera stanowi osobny poziom abstrakcji i nie jest konieczna do
budowania własnych widoków.

\section{Różnica między \texttt{transform} i \texttt{views::transform}}

Poniższe operacje mogą prowadzić do takich samych liczb, ale różnią
się \emph{czasem wykonania} i \emph{sposobem przechowywania wyniku}:

\begin{lstlisting}
std::vector<int> source{1, 2, 3};
std::vector<int> eager;
std::ranges::transform(source,
    std::back_inserter(eager),
    [](int x) { return x * x; });
// eager owns {1, 4, 9} right now

auto lazy = source | std::views::transform(
    [](int x) { return x * x; });
// lazy does not own a vector of squared values
\end{lstlisting}

Wariant pierwszy wymaga pamięci na trzy elementy i kończy obliczenia
przed kolejną instrukcją. Wariant drugi może obliczać wartości dopiero
przy dereferencji iteratorów. Po zmianie \texttt{source} kolejne
odczyty widoku \texttt{lazy} mogą odzwierciedlać nowe dane.

\section{Projektowanie API z wykorzystaniem Ranges}
\label{sec:cpp8-api}

Jeżeli funkcja tylko przegląda elementy dowolnego źródła, preferuj
parametr ograniczony koncepcją \texttt{input\_range}. Dla wielokrotnych
przejść korzystaj z \texttt{forward\_range}. Jeśli potrzebujesz
ciągłego bloku pamięci i znasz typ elementów, wygodnym wyborem jest
\texttt{std::span<T>}.

\begin{lstlisting}
template <std::ranges::input_range R>
void print_each(R&& r) {
    for (auto&& x : r) {
        // Print or consume x.
    }
}

void process_buffer(std::span<const double> values) {
    for (double value : values) {
        // Read numeric values.
    }
}
\end{lstlisting}

Zależnie od celu można zwracać widok albo kontener. Widok jest
wygodny, gdy chcemy dalej budować potok i mamy kontrolę nad czasem
życia źródła. Kontener jest często bezpieczniejszy, gdy zwracany
wynik ma istnieć niezależnie od wszystkich argumentów.

\section{Najczęstsze błędy i praktyczne reguły}
\label{sec:cpp8-pitfalls}

\begin{enumerate}
\item \textbf{Dereferencja \texttt{end()}.} Iterator końcowy jest
granicą, a nie elementem. Przed \texttt{*it} trzeba mieć
\texttt{it != end}.

\item \textbf{Mylenie kategorii iteratorów.} Lista nie obsługuje
skoków w stałym czasie. Nie można użyć dla niej algorytmu
\texttt{ranges::sort}; należy wywołać metodę \texttt{list::sort}.

\item \textbf{Traktowanie widoku jak kontenera.} \texttt{views::filter}
nie tworzy kopii wybranych elementów. Zmiany i zniszczenie źródła
mogą wpływać na wynik.

\item \textbf{Wiszące referencje w lambdach.} Zwrócony widok
z predykatem przechwytującym lokalny obiekt przez referencję może
być niebezpieczny nawet wtedy, gdy sam kontener ma odpowiednio długi
czas życia.

\item \textbf{Efekty uboczne w predykacie.} Liczba i kolejność
wywołań w leniwym potoku nie są odpowiednią podstawą działania
programu; predykaty filtrowania powinny zachować stabilną semantykę.

\item \textbf{Zakładanie stałego kosztu operacji.} Obliczenie
\texttt{distance} lub pierwszego pasującego elementu może wymagać
przejrzenia dużej części źródła.

\item \textbf{Mieszanie standardów.} \texttt{ranges::to},
\texttt{views::zip} i \texttt{views::enumerate} nie należą do C++20.

\item \textbf{Używanie nieważnych iteratorów.} Realokacja wektora
lub usuwanie elementu może unieważnić iteratory także wtedy, gdy
znajdują się wewnątrz wcześniej utworzonego widoku.
\end{enumerate}

\section{Której konstrukcji użyć?}

\begin{center}
\small
\begin{tabular}{|p{6.0cm}|p{6.5cm}|}
\hline
\textbf{Potrzeba} & \textbf{Zalecany wybór} \\
\hline
Przejście po zwykłym kontenerze & Pętla \texttt{for (auto\& x : v)} \\
\hline
Jednorazowe sortowanie albo wyszukiwanie & \texttt{std::ranges::sort}, \texttt{std::ranges::find} \\
\hline
Leniwe filtrowanie i przekształcanie & \texttt{std::views::filter}, \texttt{std::views::transform} \\
\hline
Generowanie sekwencji indeksów & \texttt{std::views::iota} \\
\hline
Odczyt lub modyfikacja ciągłego bufora & \texttt{std::span<T>} (bez własności) \\
\hline
Lekki odczyt fragmentu tekstu & \texttt{std::string\_view} (bez własności) \\
\hline
Zakres wyznaczony iteratorami & \texttt{std::ranges::subrange} \\
\hline
Wynik niezależny od źródła & Kontener, np. \texttt{std::vector} \\
\hline
Własny generator wartości & Typ modelujący \texttt{std::ranges::view} \\
\hline
Wydajny ogólny algorytm generyczny & Szablon z \texttt{std::ranges::input\_range} lub silniejszą koncepcją \\
\hline
\end{tabular}
\end{center}

\paragraph{Pytania kontrolne.}
\begin{enumerate}
\item Dlaczego iterator \texttt{end()} nie jest elementem kolekcji?
\item Co odróżnia \texttt{forward\_iterator} od \texttt{input\_iterator}?
\item Czy \texttt{std::ranges::sort} działa dla \texttt{std::list}? Dlaczego?
\item Po co \texttt{std::ranges} dopuszcza sentinel innego typu niż iterator?
\item W którym momencie \texttt{views::transform} wykonuje funkcję?
\item Dlaczego \texttt{views::filter} nie powinien zmieniać elementów, od których zależy jego predykat?
\item Czym różni się \texttt{views::all(lvalue)} od \texttt{views::all(rvalue)} dla \texttt{std::vector}?
\item Co gwarantuje, a czego nie gwarantuje \texttt{borrowed\_range}?
\item Dlaczego \texttt{ranges::find(std::vector\{\dots\}, x)} może zwrócić \texttt{dangling}?
\item Jak zebrać elementy \texttt{filter\_view} do wektora w czystym C++20?
\item Czym różni się \texttt{std::span<const T>} od \texttt{std::vector<T>}?
\item Co należy zaimplementować, aby własny obiekt był \texttt{forward\_range}?
\end{enumerate}

\paragraph{Dalsza lektura.}
Dobrym punktem odniesienia jest dokumentacja \emph{cppreference}:
\emph{Ranges library}, \emph{Iterator library}, \emph{Range adaptors},
\emph{std::ranges::borrowed\_range}, \emph{std::ranges::dangling}
oraz \emph{std::ranges::view\_interface}. Warto testować kod z opcjami kompilatora
\texttt{-std=c++20} i \texttt{-Wall -Wextra -Wpedantic},
a także korzystać z AddressSanitizer przy badaniu czasu życia danych.


\chapter[Iteratory, Ranges i Views]{Iteratory, zakresy i widoki: \textit{Ranges} oraz \textit{Views} w C++20}
\label{chap:ranges-views-cpp}

\section{Dlaczego wprowadzono Ranges?}

Klasyczne algorytmy STL operują na \emph{parze iteratorów}: pierwszym
wskazującym początek danych i drugim wyznaczającym koniec. Dzięki temu
\texttt{std::sort} lub \texttt{std::find} mogą współpracować z wieloma
kontenerami. Wymaga to jednak przekazywania obu granic i dodatkowej
ostrożności, aby pochodziły z tego samego zakresu.

\begin{lstlisting}
#include <algorithm>
#include <vector>

std::vector<int> xs{4, 1, 3, 2};
std::sort(xs.begin(), xs.end()); // STL before C++20
std::ranges::sort(xs);          // C++20
\end{lstlisting}

Nowy model pozwala przekazać \emph{zakres} (\textit{range}) jako jeden
obiekt oraz formułować wymagania wobec niego poprzez \emph{koncepcje}
(\textit{concepts}, rozdział~5). Oddzielną innowacją są \emph{widoki}
(\textit{views}): lekkie reprezentacje danych, często przekształcane
\textbf{leniwe}, to znaczy dopiero podczas ich przeglądania.

\begin{lstlisting}
#include <ranges>

std::vector<int> xs{1, 2, 3, 4, 5, 6};
auto pipeline = xs
    | std::views::filter([](int x) { return x % 2 == 0; })
    | std::views::transform([](int x) { return x * x; });

// No vector of squares has been created.
// Iteration produces 4, 16, 36.
for (int x : pipeline) { /* use x */ }
\end{lstlisting}

\paragraph{Założenia.}
W całym rozdziale przyjmujemy standard \textbf{C++20} oraz nagłówki
\texttt{<algorithm>}, \texttt{<iterator>}, \texttt{<ranges>},
\texttt{<concepts>}, \texttt{<vector>} i pozostałe nagłówki wymienione
w przykładach. W dalszych sekcjach posłużymy się semantyką przenoszenia
(rozdział~6) i lambdami (rozdział~7).

\section{Iteratory: uogólnienie wskaźników}
\label{sec:cpp8-iterators}

\subsection{Operacje \texttt{begin}, \texttt{end}, \texttt{*} i \texttt{++}}

Iterator jest obiektem umożliwiającym poruszanie się po elementach
zakresu. Wskaźnik \texttt{int*} jest przykładem iteratora, ale iterator
może także przechowywać indeks, wskaźnik do węzła listy lub stan
wewnętrzny generatora. Podstawowe operacje przedstawia program:

\begin{lstlisting}
#include <vector>

std::vector<int> v{10, 20, 30};
auto it = v.begin();     // points to 10
int a = *it;            // 10
++it;                   // points to 20
int b = *it;            // 20
bool ok = it != v.end();
\end{lstlisting}

Iterator \texttt{end()} zwykle oznacza pozycję \emph{za ostatnim}
elementem; nie wolno go dereferencjonować. Dzięki temu pusty zakres
spełnia \texttt{begin()==end()}, a przechodzenie po elementach ma
jednolity schemat:

\begin{lstlisting}
for (auto it = v.begin(); it != v.end(); ++it) {
    // *it is a valid element
}
\end{lstlisting}

\paragraph{Pętla zakresowa.}
Zapis \texttt{for (auto\& x : v)} jest wygodniejszą postacią iteracji,
ale w dalszym ciągu wykorzystuje operacje początku, końca,
porównania, zwiększania i dereferencji. Aby \emph{modyfikować} elementy,
stosujemy \texttt{auto\&}; aby uniknąć kopiowania bez modyfikowania,
\texttt{const auto\&}; dla niewielkich wartości liczbowych często
wystarczy \texttt{auto}.

\subsection{Iteratory stałe i nieważność iteratorów}

\begin{lstlisting}
std::vector<int> v{1, 2, 3};
auto m = v.begin();   // can modify *m
auto c = v.cbegin();  // cannot modify *c
*m = 10;
// *c = 20;         // compilation error
\end{lstlisting}

\texttt{const\_iterator} chroni elementy przed modyfikacją poprzez
ten iterator; nie oznacza, że cały kontener jest niezmienny. Równie
ważna jest \emph{nieważność iteratorów} (\textit{iterator invalidation}).
Na przykład ponowna alokacja pamięci przez \texttt{std::vector::push\_back}
może unieważnić \textbf{wszystkie} iteratory, wskaźniki i referencje do
jego elementów.

\begin{lstlisting}
std::vector<int> v{1, 2, 3};
auto it = v.begin();
v.push_back(4);       // may reallocate storage
// *it;              // unsafe if reallocation happened
it = v.begin();      // obtain a fresh iterator
\end{lstlisting}

Reguły zależą od kontenera i operacji. W \texttt{std::list} wstawienie
elementu nie unieważnia iteratorów do pozostałych elementów, ale nie
oznacza to, że można dalej używać iteratora do elementu usuniętego.

\section{Kategorie iteratorów i ich wymagania}
\label{sec:cpp8-iterator-categories}

Zestaw koncepcji C++20 opisuje operacje i własności iteratorów, a nie
jedynie nazwy klas. Najważniejsze kategorie zestawia tabela.

\begin{center}
\small
\begin{tabular}{|p{5.2cm}|p{7.3cm}|}
\hline
\textbf{Koncepcja} & \textbf{Możliwości i przykłady} \\
\hline
\texttt{input\_iterator} & Odczytywanie kolejnych wartości; przeglądanie może być jednokrotne (np. strumień wejściowy). \\
\hline
\texttt{output\_iterator} & Zapisywanie kolejnych wartości do celu (np. \texttt{back\_insert\_iterator}). \\
\hline
\texttt{forward\_iterator} & Wielokrotne przejścia, porównywanie iteratorów, krok do przodu. \\
\hline
\texttt{bidirectional\_iterator} & Dodatkowo krok wstecz, czyli \texttt{--it}; np. \texttt{list}. \\
\hline
\texttt{random\_access\_iterator} & Indeksowanie i skoki \texttt{it+n} w stałym czasie; np. \texttt{deque}. \\
\hline
\texttt{contiguous\_iterator} & Elementy leżą kolejno w pamięci; np. iterator \texttt{vector<int>} i surowy wskaźnik. \\
\hline
\end{tabular}
\end{center}

Hierarchia wymagań dla iteratorów odczytujących jest w przybliżeniu
następująca (każda kolejna koncepcja jest silniejsza):

\begin{center}
\texttt{input $\Rightarrow$ forward $\Rightarrow$ bidirectional}
\\[2pt]
\texttt{$\Rightarrow$ random\_access $\Rightarrow$ contiguous}.
\end{center}

Koncepcja \texttt{output\_iterator} dotyczy osobnej rodziny wymagań.
Nowe koncepcje nie są w każdym szczególe identyczne ze starymi
kategoriami \emph{LegacyIterator}. Na przykład model
\texttt{forward\_iterator} nie wymaga zwracania referencji
przy dereferencji.

\begin{lstlisting}
#include <concepts>
#include <iterator>
#include <list>
#include <vector>

static_assert(std::contiguous_iterator<int*>);
static_assert(std::random_access_iterator<
    std::vector<int>::iterator>);
static_assert(std::bidirectional_iterator<
    std::list<int>::iterator>);
static_assert(!std::random_access_iterator<
    std::list<int>::iterator>);
\end{lstlisting}

\paragraph{Dlaczego to ważne?}
\texttt{std::ranges::sort} wymaga iteratorów o dostępie swobodnym,
więc można użyć go do \texttt{std::vector}, ale \textbf{nie} do
\texttt{std::list}. Listę sortuje się metodą \texttt{list::sort()}.
Wybór algorytmu jest zatem ograniczony kategorią iteratora.

\section{Własny iterator i własny zakres}
\label{sec:cpp8-own-iterator}

Zbudujemy iterator generujący kolejne liczby całkowite. Nie będzie on
przechowywać ani wskaźnika, ani kontenera: całą potrzebną informacją
jest bieżąca wartość. Zdefiniujemy \emph{iterator do przodu}
(\texttt{forward\_iterator}).

\begin{lstlisting}
#include <concepts>
#include <cstddef>
#include <iterator>
#include <ranges>

struct CountingIterator {
    using value_type = int;
    using difference_type = std::ptrdiff_t;
    using iterator_concept = std::forward_iterator_tag;

    int current = 0;

    int operator*() const { return current; }

    CountingIterator& operator++() {
        ++current;
        return *this;
    }
    CountingIterator operator++(int) {
        auto old = *this;
        ++*this;
        return old;
    }
    bool operator==(const CountingIterator&) const = default;
};

struct CountingRange {
    int first = 0;
    int last = 0;

    CountingIterator begin() const { return {first}; }
    CountingIterator end() const { return {last}; }
};

static_assert(std::forward_iterator<CountingIterator>);
static_assert(std::ranges::forward_range<CountingRange>);
\end{lstlisting}

Teraz kod \texttt{for (int i : CountingRange\{3, 7\})} wypisze
\texttt{3, 4, 5, 6}. Funkcje \texttt{begin()} i \texttt{end()}
wyznaczają \emph{półotwarty} zakres \([3,7)\). Prawidłowe użycie
zakłada \texttt{first <= last}; dla \texttt{first > last} nasz
minimalny iterator nie zbiegnie do wartości końcowej.

\paragraph{Co trzeba zagwarantować?}
Samo przejście \texttt{static\_assert} nie dowodzi wszystkich
semantycznych własności iteratora. Dla \texttt{forward\_iterator}
ważna jest m.in. własność \emph{multi-pass}: kopiowanie iteratora
i przejście jedną kopią nie powinno zmieniać wartości zwracanych
przez drugą. Nasz iterator ją spełnia, ponieważ każda kopia ma własne
pole \texttt{current}.

\section{Zwykłe \textit{range} i \textit{sentinel}}
\label{sec:cpp8-range-sentinel}

\emph{Zakres} jest obiektem, dla którego można otrzymać początek
\texttt{std::ranges::begin(r)} i koniec \texttt{std::ranges::end(r)}.
Typ wyniku obu wywołań \textbf{nie musi być identyczny}. Obiekt
wyznaczający koniec to \emph{sentinel}; musi dać się porównać z
iteratorem, ale sam nie musi obsługiwać jego pozostałych operacji.

\begin{lstlisting}
#include <ranges>
#include <vector>

std::vector<int> xs{1, 2, 3};
auto first = std::ranges::begin(xs);
auto last = std::ranges::end(xs);
while (first != last) { ++first; }
\end{lstlisting}

W tym przypadku obie granice mają ten sam typ: zakres spełnia
koncepcję \texttt{common\_range}. W innych zakresach koniec może
wyznaczać osobny obiekt sentinela. Odpowiednia koncepcja
\texttt{std::sentinel\_for} określa warunki porównywania granic:

\begin{lstlisting}
static_assert(std::ranges::common_range<decltype(xs)>);
static_assert(std::sentinel_for<
    decltype(last), decltype(first)>);
\end{lstlisting}

Przykładem osobnego typu końca jest \texttt{std::default\_sentinel\_t}.

\paragraph{Rozmiar a liczba kroków.}
Operacja \texttt{std::ranges::size(r)} jest dostępna dla zakresów
spełniających \texttt{sized\_range}. Nie każdy zakres ma tanią operację
wyznaczenia rozmiaru. Funkcja \texttt{distance(r)} z przestrzeni
\texttt{std::ranges} może wymagać
\emph{przejścia przez wszystkie elementy}. Dla niektórych
jednokrotnych źródeł danych przejście oznacza również ich zużycie.

\section{Koncepcje zakresów w praktyce}

Własności zakresów sprawdzamy za pomocą koncepcji z rozdziału~5.
Najważniejsze z nich to:
\begin{itemize}
\item \texttt{input\_range} i \texttt{forward\_range},
\item \texttt{bidirectional\_range} i \texttt{random\_access\_range},
\item \texttt{contiguous\_range}, \texttt{common\_range} i \texttt{sized\_range}.
\end{itemize}

\begin{lstlisting}
#include <ranges>
#include <vector>

template <std::ranges::input_range R>
long long sum_integers(R&& range) {
    long long result = 0;
    for (auto value : range) result += value;
    return result;
}

std::vector<int> v{1, 2, 3};
long long s = sum_integers(v); // 6
\end{lstlisting}

Parametr \texttt{R\&\&} jest tutaj \emph{referencją przekazującą}
(\textit{forwarding reference}), a nie żądaniem przejęcia własności.
Funkcja obsługuje także źródła jednokrotne, gdyż przechodzi zakres tylko
raz. Dla naprawdę uniwersalnej funkcji trzeba osobno wyrazić warunki
na typ elementu i wykonywaną na nim operację dodawania.

\section{Algorytmy \texttt{std::ranges} i projekcje}
\label{sec:cpp8-algorithms}

Funkcje z rodziny \texttt{std::ranges}, np. \texttt{sort},
\texttt{find}, \texttt{copy} i \texttt{transform}, przyjmują
bezpośrednio zakres albo parę iterator--sentinel.
Ich wykonanie jest na ogół \emph{natychmiastowe} (w odróżnieniu od
samych widoków \texttt{views::transform}).

\begin{lstlisting}
#include <algorithm>
#include <iterator>
#include <vector>

std::vector<int> v{4, 2, 3, 1};
std::ranges::sort(v);  // modifies v immediately

auto pos = std::ranges::find(v, 3);
if (pos != v.end()) { *pos = 30; }

std::vector<int> out;
std::ranges::copy(v, std::back_inserter(out));
\end{lstlisting}

\subsection{Projekcja: sortowanie po wskazanym polu}

Algorytmy Ranges często posiadają argument \emph{projekcji}: funkcję
wywoływaną na elemencie przed porównaniem albo wyszukiwaniem. Dzięki
temu zamiast pisać własny komparator, można wskazać pole struktury.

\begin{lstlisting}
#include <string>
#include <vector>

struct Person {
    std::string name;
    int age;
};

std::vector<Person> people{
    {"Ada", 40}, {"Ewa", 25}, {"Jan", 32}
};

std::ranges::sort(people, {}, &Person::age);
// ascending ages: 25, 32, 40

auto it = std::ranges::find(people, 32, &Person::age);
// it->name == "Jan"
\end{lstlisting}

W wywołaniu \texttt{sort(people, \{\}, \&Person::age)} puste
nawiasy oznaczają domyślny komparator \texttt{std::ranges::less},
a wskaźnik do składowej to projekcja. Standard używa mechanizmu
\texttt{std::invoke}, omówionego w rozdziale~7.

\section{Widoki: \textit{view} nie jest synonimem kontenera}
\label{sec:cpp8-views}

\emph{Widok} jest zakresem przystosowanym do taniego tworzenia i
kopiowania lub przenoszenia \emph{samego obiektu widoku}. Nie musi mieć
własnych danych: może przechowywać referencję do kontenera, mały zestaw
iteratorów albo parametry generowania wartości.

\begin{lstlisting}
#include <ranges>
#include <vector>

std::vector<int> xs{10, 20, 30, 40};
auto first_two = xs | std::views::take(2);
// first_two refers to the original xs
\end{lstlisting}

Stworzenie widoku nie jest równoznaczne ze skopiowaniem elementów do
nowego kontenera. Odwołanie do danych źródłowych wymaga, by pozostawały
one ważne przez cały czas korzystania z widoku. Szczegółowe reguły
własności omówimy dalej.

\paragraph{Dwa znaczenia nazwy \texttt{views}.}
\texttt{std::views} jest aliasem przestrzeni nazw z obiektami
adaptującymi zakresy. \texttt{std::ranges::view} to natomiast
\emph{koncepcja} spełniana przez dany typ. Nie należy mylić ich
z pojęciem \texttt{std::ranges::range}.

\section{Leniwe potoki \texttt{filter}, \texttt{transform}, \texttt{take}}
\label{sec:cpp8-pipelines}

Najczęściej stosowane widoki łączymy za pomocą operatora \texttt{|}.
Adaptory przekazują sobie zakresy bez potrzeby budowania pośrednich
wektorów:

\begin{lstlisting}
std::vector<int> xs{1,2,3,4,5,6,7,8,9,10};

auto pipeline = xs
    | std::views::filter([](int x) {
          return x % 2 == 0;
      })
    | std::views::transform([](int x) {
          return x * x;
      })
    | std::views::take(3);

// Iteration produces 4, 16, 36.
\end{lstlisting}

\texttt{filter} pomija elementy niespełniające predykatu,
\texttt{transform} wyznacza nową wartość przy dereferencji,
a \texttt{take(3)} ogranicza liczbę wyników. Algorytm
\texttt{std::ranges::transform} działa od razu; natomiast
utworzenie widoku \texttt{views::transform} nie przelicza
wszystkich wyników.

\subsection{Kiedy następują obliczenia?}

\begin{lstlisting}
int calls = 0;
auto squares = xs | std::views::transform([&](int x) {
    ++calls;
    return x * x;
});
// calls == 0 here
int first = *squares.begin();
// calls has increased (transform was evaluated)
\end{lstlisting}

W tym przykładzie pokazujemy moment obliczenia, a nie zalecany styl
projektowania: \textbf{predykaty i transformacje powinny zasadniczo
unikać efektów ubocznych}. Ponowne przejście widoku, ponowna
dereferencja lub implementacyjnie dopuszczone buforowanie może
powodować inne liczby wywołań, niż sugerowałaby pojedyncza pętla.
Nie należy projektować logiki programu na podstawie rzekomej liczby
wywołań funkcji przez widok.

\paragraph{Ważny szczegół \texttt{filter\_view}.}
Dla pewnych kategorii zakresu \texttt{filter\_view} może zapamiętywać
pierwszy znaleziony element. Ponadto modyfikowanie elementu tak,
aby przestał spełniać predykat filtrowania, narusza wymagania
\texttt{filter\_view} i może prowadzić do niezdefiniowanego
zachowania. Predykat powinien być stabilny wobec przebiegu iteracji.

\section{Pozostałe przydatne adaptory C++20}
\label{sec:cpp8-adaptors}

\subsection{\texttt{iota}: wartości wyznaczane bez tablicy}

\begin{lstlisting}
auto ids = std::views::iota(1, 5);
// Iteration produces 1, 2, 3, 4.

// This is an unbounded range, limited by take(5):
auto first_five = std::views::iota(1)
                | std::views::take(5);
// Produces 1, 2, 3, 4, 5.
\end{lstlisting}

\texttt{iota} przydaje się do indeksowania lub generowania sekwencji.
Dla nieograniczonego widoku trzeba pamiętać o warunku zakończenia;
ponadto operacje arytmetyczne nie mogą przepełniać typu całkowitego
ze znakiem.

\subsection{\texttt{drop}, \texttt{take} i \texttt{reverse}}

\begin{lstlisting}
std::vector<int> data{1, 2, 3, 4, 5};
auto middle = data | std::views::drop(1)
                   | std::views::take(3);
// 2, 3, 4

auto backwards = data | std::views::reverse;
// 5, 4, 3, 2, 1
\end{lstlisting}

\texttt{reverse} wymaga co najmniej zakresu dwukierunkowego;
\texttt{take} i \texttt{drop} są przydatne także dla sekwencji
generowanych leniwie. Dla niektórych widoków koszt \texttt{begin()}
lub dojścia do pierwszego wyniku może być większy niż stały.

\subsection{\texttt{keys}, \texttt{values} i \texttt{split}}

\begin{lstlisting}
std::vector<std::pair<int, std::string>> pairs{
    {1, "one"}, {2, "two"}
};
for (int key : pairs | std::views::keys) { /* 1, 2 */ }
for (const auto& value : pairs | std::views::values) {
    /* "one", "two" */
}

std::string csv = "aa,b,cc";
for (auto part : csv | std::views::split(',')) {
    std::string piece(part.begin(), part.end());
    // Copy each part into an owning string if needed.
}
\end{lstlisting}

\texttt{split} nie tworzy automatycznie \texttt{std::string} dla
każdego fragmentu; otrzymujemy zakresy wskazujące fragmenty źródła.
Zachowując wynik na później, trzeba zapewnić odpowiedni czas życia
źródła. Nie należy też zakładać, że splitowanie tekstu CSV obsługuje
cudzysłowy i escapowanie pól zgodnie z pełnym formatem CSV.

\paragraph{Uwaga na wersję standardu.}
\texttt{std::views::zip} i \texttt{std::views::enumerate} są
konstrukcjami \textbf{C++23}, podobnie jak \texttt{std::ranges::to}.
Nie są dostępne w czystym C++20. W standardzie C++20 wynik można zebrać do kontenera np.
algorytmem \texttt{ranges::copy} z iteratorem wstawiającym
\texttt{std::back\_inserter}.

\section{Materializacja wyniku i koszt przetwarzania}
\label{sec:cpp8-materialize}

Jeśli wynik ma być niezależny od źródła i przechowywany na później,
możemy zamienić potok w kontener. W C++20 nie ma standardowego
\texttt{std::ranges::to<std::vector>()}.

\begin{lstlisting}
#include <algorithm>
#include <iterator>
#include <vector>

std::vector<int> data{1,2,3,4,5,6};
auto even = data | std::views::filter([](int x) {
    return x % 2 == 0;
});

std::vector<int> result;
std::ranges::copy(even, std::back_inserter(result));
// result == {2, 4, 6}
\end{lstlisting}

Taka materializacja wymaga pamięci proporcjonalnej do wielkości wyniku.
Zaletą potoku jest unikanie pośrednich alokacji i często możliwość
\emph{łączenia} transformacji w jednym przebiegu. Nie jest to jednak
gwarancja, że widoki zawsze będą szybsze: koszty iteracji, predykatów,
przepływu sterowania i optymalizacji zależą od konkretnego kodu i
kompilatora. W zadaniach krytycznych wykonuj pomiary.

\section{Własność danych, czas życia i \texttt{borrowed\_range}}
\label{sec:cpp8-lifetime}

To najważniejsza część praktycznego korzystania z Ranges. Widoki
\textbf{nie zawsze posiadają} bazowe dane. W przypadku
\texttt{std::views::all} przeprowadzamy istotne rozróżnienie:

\begin{lstlisting}
std::vector<int> source{2, 4, 6};
auto borrowed = std::views::all(source);
// Refers to source: source must stay alive.

auto owned = std::views::all(std::vector<int>{3, 5, 7});
// In C++20, this can be an owning_view.
// The temporary vector is moved into the view.
\end{lstlisting}

Nie oznacza to, że \emph{każdy} widok utworzony z wartości tymczasowej
przejmuje własność, ani że wszystkie widoki zachowują się tak samo.
Różne adaptory przechowują różne rodzaje obiektów bazowych.

\subsection{Wskaźniki wiszące wewnątrz potoku}

\begin{lstlisting}
auto make_bad_view() {
    std::vector<int> local{1, 2, 3};
    return local | std::views::filter([](int x) {
        return x > 1;
    });
}
// Using the returned view would access destroyed local.
\end{lstlisting}

Ponieważ \texttt{local} jest \emph{lvalue}, adapter przechowuje
nieposiadający danych widok do lokalnego kontenera. Po wyjściu z
funkcji \texttt{local} zostaje zniszczony, a zwrócony widok zawiera
nieaktualne odwołanie. Bezpieczne wyjścia to utworzenie widoku z
przekazanego na własność kontenera (np. obiektu tymczasowego lub
\texttt{std::move(local)}) albo zwrócenie zwykłego kontenera wynikowego.

\begin{lstlisting}
auto make_owning_view() {
    std::vector<int> local{1, 2, 3};
    return std::views::all(std::move(local))
         | std::views::filter([](int x) { return x > 1; });
}
// The returned view owns its vector through owning_view.
\end{lstlisting}

\subsection{Co naprawdę mówi \texttt{borrowed\_range}?}

Koncepcja \texttt{borrowed\_range<R>} oznacza w uproszczeniu, że
iteratory otrzymane z zakresu typu \texttt{R} mogą pozostać ważne
po zniszczeniu \emph{samego obiektu zakresu} (o ile dane, do których
odsyłają, nadal istnieją). Nie jest to synonim ``widok nie posiada
danych'' ani ``wszystkie wskaźniki są zawsze bezpieczne''.

\begin{lstlisting}
#include <concepts>
#include <ranges>
#include <vector>

static_assert(std::ranges::borrowed_range<
    std::ranges::ref_view<std::vector<int>>>);
static_assert(!std::ranges::borrowed_range<
    std::vector<int>>);
\end{lstlisting}

\subsection{\texttt{std::ranges::dangling} jako zabezpieczenie}

Algorytm zwracający iterator do \emph{tymczasowego} kontenera nie może
bezpiecznie zwrócić rzeczywistego iteratora do zniszczonych danych.
Dlatego pewne algorytmy Ranges zwracają specjalny typ
\texttt{std::ranges::dangling}.

\begin{lstlisting}
#include <type_traits>

using Result = decltype(std::ranges::find(
    std::vector<int>{1, 2, 3}, 2));
static_assert(std::same_as<Result, std::ranges::dangling>);
\end{lstlisting}

Nie należy jednak oczekiwać, że mechanizm ten wykryje wszystkie
problemy z czasem życia. Jeżeli zapisujemy nieposiadający danych
widok do lokalnej zmiennej, kompilator może nie zauważyć, że docelowy
kontener umrze wcześniej niż widok.

\section{\texttt{std::span}, \texttt{std::string\_view} i \texttt{subrange}}
\label{sec:cpp8-span}

\texttt{std::span<T>} opisuje ciągły fragment pamięci jako wskaźnik
i rozmiar (lub rozmiar znany w czasie kompilacji). Jest to lekka,
\textbf{nieposiadająca danych} alternatywa dla pary
\texttt{(T*, size)}.

\begin{lstlisting}
#include <span>

void normalize(std::span<double> xs) {
    for (double& x : xs) x /= 100.0;
}

std::vector<double> v{100, 200, 300};
normalize(v);  // v becomes {1, 2, 3}
\end{lstlisting}

\texttt{std::span<const T>} umożliwia odczyt bez modyfikowania
elementów, a \texttt{std::span<T>} pozwala modyfikować elementy, mimo
że sam widok nie posiada pamięci.

\begin{lstlisting}
int a[]{1, 2, 3, 4};
std::span<int> all{a};
auto head = all.first(2);
head[1] = 20; // a[1] == 20
\end{lstlisting}

\texttt{std::string\_view} spełnia podobną rolę dla tekstu; nie
alokuje pamięci na znaki (oraz nie musi kończyć się bajtem zero).
Trzeba uważać, aby nie zapisać go po zniszczeniu źródłowego
\texttt{std::string}.

\begin{lstlisting}
#include <string_view>

std::string text = "abcdef";
std::string_view view{text};
std::string_view prefix = view.substr(0, 3); // "abc"
// Do not use view after text is destroyed or invalidated.
\end{lstlisting}

\texttt{std::ranges::subrange<I, S>} grupuje iterator \texttt{I}
i sentinel \texttt{S} w zakres bez kopiowania elementów. Także tutaj
obowiązuje odpowiedzialność za czas życia wskazywanych danych.

\begin{lstlisting}
auto sub = std::ranges::subrange(v.begin(), v.end());
// sub refers to elements of v; it does not own v.
\end{lstlisting}

\paragraph{Ostrożnie z realokacją.}
\texttt{std::span} do elementów wektora, \texttt{string\_view} do
znaków łańcucha oraz \texttt{subrange} z iteratorów do kontenera mogą
stać się niepoprawne po operacji zmieniającej miejsce przechowywania
danych. Przenoszenie lub realokacja nie są automatycznie śledzone.

\section{Własny widok za pomocą \texttt{view\_interface}}
\label{sec:cpp8-custom-view}

Na końcu zbudujemy \emph{własny widok} generujący kwadraty liczb
całkowitych od zera do \(n-1\). Użyjemy
\texttt{std::ranges::view\_interface}, które dostarcza wygodnych
operacji opartych na \texttt{begin()} i \texttt{end()}. W klasie
pozostają wyłącznie parametry generatora, więc nie ma alokacji pamięci
na elementy sekwencji.

\begin{lstlisting}
#include <concepts>
#include <cstddef>
#include <iterator>
#include <ranges>

struct SquaresView : std::ranges::view_interface<SquaresView> {
    struct iterator {
        using value_type = int;
        using difference_type = std::ptrdiff_t;
        using iterator_concept = std::forward_iterator_tag;

        int index = 0;
        int operator*() const { return index * index; }
        iterator& operator++() {
            ++index;
            return *this;
        }
        iterator operator++(int) {
            auto old = *this;
            ++*this;
            return old;
        }
        bool operator==(const iterator&) const = default;
    };

    int n = 0;
    SquaresView() = default;
    explicit SquaresView(int count) : n(count) {}

    iterator begin() const { return {0}; }
    iterator end() const { return {n}; }
};

static_assert(std::ranges::view<SquaresView>);
static_assert(std::ranges::forward_range<SquaresView>);
\end{lstlisting}

\begin{lstlisting}
SquaresView squares{6};
for (int x : squares) { /* 0, 1, 4, 9, 16, 25 */ }

// Our view works with standard adapters:
auto even_squares = squares | std::views::filter(
    [](int x) { return x % 2 == 0; });
// Iteration yields 0, 4, 16.
\end{lstlisting}

Konstruktor z argumentem \texttt{count} powinien otrzymać wartość
nieujemną. Ponieważ liczymy \texttt{index * index} w typie
\texttt{int}, należy ponadto ograniczyć rozmiar zakresu, aby uniknąć
przepełnienia. W rozwiązaniu produkcyjnym należałoby te warunki
weryfikować oraz dopracować obsługę szerszych typów całkowitych.

\paragraph{Widok a adapter.}
\texttt{SquaresView} jest własnym \emph{widokiem}, który można
przekazać do standardowych adapterów. Nie jest jeszcze własnym
\emph{range adaptor closure}, czyli obiektem pozwalającym zapisać
\texttt{some\_range | squares\_adaptor}. Implementacja takiego
adaptera stanowi osobny poziom abstrakcji i nie jest konieczna do
budowania własnych widoków.

\section{Różnica między \texttt{transform} i \texttt{views::transform}}

Poniższe operacje mogą prowadzić do takich samych liczb, ale różnią
się \emph{czasem wykonania} i \emph{sposobem przechowywania wyniku}:

\begin{lstlisting}
std::vector<int> source{1, 2, 3};
std::vector<int> eager;
std::ranges::transform(source,
    std::back_inserter(eager),
    [](int x) { return x * x; });
// eager owns {1, 4, 9} right now

auto lazy = source | std::views::transform(
    [](int x) { return x * x; });
// lazy does not own a vector of squared values
\end{lstlisting}

Wariant pierwszy wymaga pamięci na trzy elementy i kończy obliczenia
przed kolejną instrukcją. Wariant drugi może obliczać wartości dopiero
przy dereferencji iteratorów. Po zmianie \texttt{source} kolejne
odczyty widoku \texttt{lazy} mogą odzwierciedlać nowe dane.

\section{Projektowanie API z wykorzystaniem Ranges}
\label{sec:cpp8-api}

Jeżeli funkcja tylko przegląda elementy dowolnego źródła, preferuj
parametr ograniczony koncepcją \texttt{input\_range}. Dla wielokrotnych
przejść korzystaj z \texttt{forward\_range}. Jeśli potrzebujesz
ciągłego bloku pamięci i znasz typ elementów, wygodnym wyborem jest
\texttt{std::span<T>}.

\begin{lstlisting}
template <std::ranges::input_range R>
void print_each(R&& r) {
    for (auto&& x : r) {
        // Print or consume x.
    }
}

void process_buffer(std::span<const double> values) {
    for (double value : values) {
        // Read numeric values.
    }
}
\end{lstlisting}

Zależnie od celu można zwracać widok albo kontener. Widok jest
wygodny, gdy chcemy dalej budować potok i mamy kontrolę nad czasem
życia źródła. Kontener jest często bezpieczniejszy, gdy zwracany
wynik ma istnieć niezależnie od wszystkich argumentów.

\section{Najczęstsze błędy i praktyczne reguły}
\label{sec:cpp8-pitfalls}

\begin{enumerate}
\item \textbf{Dereferencja \texttt{end()}.} Iterator końcowy jest
granicą, a nie elementem. Przed \texttt{*it} trzeba mieć
\texttt{it != end}.

\item \textbf{Mylenie kategorii iteratorów.} Lista nie obsługuje
skoków w stałym czasie. Nie można użyć dla niej algorytmu
\texttt{ranges::sort}; należy wywołać metodę \texttt{list::sort}.

\item \textbf{Traktowanie widoku jak kontenera.} \texttt{views::filter}
nie tworzy kopii wybranych elementów. Zmiany i zniszczenie źródła
mogą wpływać na wynik.

\item \textbf{Wiszące referencje w lambdach.} Zwrócony widok
z predykatem przechwytującym lokalny obiekt przez referencję może
być niebezpieczny nawet wtedy, gdy sam kontener ma odpowiednio długi
czas życia.

\item \textbf{Efekty uboczne w predykacie.} Liczba i kolejność
wywołań w leniwym potoku nie są odpowiednią podstawą działania
programu; predykaty filtrowania powinny zachować stabilną semantykę.

\item \textbf{Zakładanie stałego kosztu operacji.} Obliczenie
\texttt{distance} lub pierwszego pasującego elementu może wymagać
przejrzenia dużej części źródła.

\item \textbf{Mieszanie standardów.} \texttt{ranges::to},
\texttt{views::zip} i \texttt{views::enumerate} nie należą do C++20.

\item \textbf{Używanie nieważnych iteratorów.} Realokacja wektora
lub usuwanie elementu może unieważnić iteratory także wtedy, gdy
znajdują się wewnątrz wcześniej utworzonego widoku.
\end{enumerate}

\section{Której konstrukcji użyć?}

\begin{center}
\small
\begin{tabular}{|p{6.0cm}|p{6.5cm}|}
\hline
\textbf{Potrzeba} & \textbf{Zalecany wybór} \\
\hline
Przejście po zwykłym kontenerze & Pętla \texttt{for (auto\& x : v)} \\
\hline
Jednorazowe sortowanie albo wyszukiwanie & \texttt{std::ranges::sort}, \texttt{std::ranges::find} \\
\hline
Leniwe filtrowanie i przekształcanie & \texttt{std::views::filter}, \texttt{std::views::transform} \\
\hline
Generowanie sekwencji indeksów & \texttt{std::views::iota} \\
\hline
Odczyt lub modyfikacja ciągłego bufora & \texttt{std::span<T>} (bez własności) \\
\hline
Lekki odczyt fragmentu tekstu & \texttt{std::string\_view} (bez własności) \\
\hline
Zakres wyznaczony iteratorami & \texttt{std::ranges::subrange} \\
\hline
Wynik niezależny od źródła & Kontener, np. \texttt{std::vector} \\
\hline
Własny generator wartości & Typ modelujący \texttt{std::ranges::view} \\
\hline
Wydajny ogólny algorytm generyczny & Szablon z \texttt{std::ranges::input\_range} lub silniejszą koncepcją \\
\hline
\end{tabular}
\end{center}

\paragraph{Pytania kontrolne.}
\begin{enumerate}
\item Dlaczego iterator \texttt{end()} nie jest elementem kolekcji?
\item Co odróżnia \texttt{forward\_iterator} od \texttt{input\_iterator}?
\item Czy \texttt{std::ranges::sort} działa dla \texttt{std::list}? Dlaczego?
\item Po co \texttt{std::ranges} dopuszcza sentinel innego typu niż iterator?
\item W którym momencie \texttt{views::transform} wykonuje funkcję?
\item Dlaczego \texttt{views::filter} nie powinien zmieniać elementów, od których zależy jego predykat?
\item Czym różni się \texttt{views::all(lvalue)} od \texttt{views::all(rvalue)} dla \texttt{std::vector}?
\item Co gwarantuje, a czego nie gwarantuje \texttt{borrowed\_range}?
\item Dlaczego \texttt{ranges::find(std::vector\{\dots\}, x)} może zwrócić \texttt{dangling}?
\item Jak zebrać elementy \texttt{filter\_view} do wektora w czystym C++20?
\item Czym różni się \texttt{std::span<const T>} od \texttt{std::vector<T>}?
\item Co należy zaimplementować, aby własny obiekt był \texttt{forward\_range}?
\end{enumerate}

\paragraph{Dalsza lektura.}
Dobrym punktem odniesienia jest dokumentacja \emph{cppreference}:
\emph{Ranges library}, \emph{Iterator library}, \emph{Range adaptors},
\emph{std::ranges::borrowed\_range}, \emph{std::ranges::dangling}
oraz \emph{std::ranges::view\_interface}. Warto testować kod z opcjami kompilatora
\texttt{-std=c++20} i \texttt{-Wall -Wextra -Wpedantic},
a także korzystać z AddressSanitizer przy badaniu czasu życia danych.
